% !TeX root = ../../Book1.tex
\chapter{Sobolev 函数的逼近、局部化与延拓}
\label{b1:ch:22}
\BookChapterNumber{b1:ch:22}
\begin{chapterabstract}
本章研究如何把 Sobolev 函数换成便于计算的函数，同时保留弱导数与范数控制。内部紧支集允许补零，并在有限指数下允许磨光；逐层分解使任意开集上的局部磨光成为全域光滑逼近。无穷指数下卷积仍有范数界，却一般不在 Sobolev 范数中收敛。要求逼近函数具有紧支集，则产生零边界空间，并使边界几何进入问题。最后通过半空间反射、Lipschitz 图表及单位分解构造有界延拓，再用截断推导一个基本的局部能量估计。
\end{chapterabstract}

\begin{learninggoals}
\begin{enumerate}
\item\label{b1:goal:22-approximation}在 \(1\le p<\infty\) 的范围内证明内部紧支集与全空间的光滑逼近，掌握 Meyers--Serrin 定理中磨光尺度、支集和可和误差的配合，并辨认 \(p=\infty\) 时强逼近失效的原因。
\item\label{b1:goal:22-zero-boundary}理解 \(W_0^{1,p}\) 的闭包定义，区分补零、零端点值和内部光滑逼近，并辨认无穷指数的额外障碍。
\item\label{b1:goal:22-extension}通过沿线绝对连续性证明半空间反射，再把反射放入 Lipschitz 边界图表，构造固定的有界线性延拓。
\item\label{b1:goal:22-localization}追踪截断导数对估计常数的影响，用紧支集逼近合法化 Sobolev 测试函数，并推导局部梯度能量界。
\end{enumerate}
\end{learninggoals}

\chapref{b1:ch:21}的完备性、弱乘积公式、沿坐标线绝对连续性与双 Lipschitz 坐标变换，是本章的直接前置。光滑单位分解与磨光函数来自\chapref{b1:ch:20}。证明安排先处理开集内部，再处理边界；这样可以看清哪些结论对任意开集成立，哪些结论必须使用边界图表。

阅读时应把三种要求分开：在域内光滑、具有内部紧支集、能够延拓到边界之外。常数函数在有界区间上已经光滑，却不能由内部紧支集函数依一阶 Sobolev 范数逼近；有界 Lipschitz 区域上的函数可以受控延拓，却仍可能在 \(W^{1,\infty}\) 范数下无法磨光逼近。这些差别并不是符号细节，而会决定后续边界条件的含义。

Meyers 与 Serrin 的《\(H=W\)》\cite{MeyersSerrin1964HW}是全域光滑稠密性的原始文献，篇幅很短，适合在读完\secref{b1:sec:2201}后对照。需要把本章放回 Sobolev 空间的系统理论时，可伴读 Adams 与 Fournier 的《Sobolev Spaces》\cite{Adams2003Sobolev}及 Brézis 的《Functional Analysis, Sobolev Spaces and Partial Differential Equations》\cite{Brezis2011Functional}。延拓中的几何组织另在\secref{b1:sec:2203}说明具体参考范围。

% !TeX root = ../../Book1.tex
\section{光滑逼近}
\label{b1:sec:2201}

\chapref{b1:ch:21}已用局部卷积建立 Sobolev 函数的运算法则。现在要把这种局部逼近接成一个在全域范数下有效的光滑函数。困难不在卷积本身，而在于一个固定半径的核可能越过边界，许多个局部误差又未必能够相加。本节为不同位置选择不同的磨光尺度，同时控制支集和误差。

Meyers 与 Serrin 的短文《\(H=W\)》\cite[pp. 1055--1056]{MeyersSerrin1964HW}证明了任意开集上的光滑稠密性。原文的 \(H^{k,p}\) 指域内光滑函数在 \(W^{k,p}\) 中的闭包，不是本书专用于 \(W^{k,2}\) 的记号 \(H^k\)。以下直接陈述稠密性，不另外引入同名空间。

% 实读来源：MeyersSerrin1964HW，PNAS 51(6) (1964), pp.1055--1056，
% DOI 10.1073/pnas.51.6.1055，PMCID PMC300210；Europe PMC正式扫描副本两页均已阅读。
% p.1055：任意开集、复值、k>=0、1<=p<infty，原文用各导数Lp范数之和；
% p.1056：局部Sobolev函数的全局小误差逼近、环层单位分解及逐层卷积证明。
% 本节用已证完备性处理可和误差级数，补明紧支集补零、支集膨胀与局部有限性；
% 原文末段关于光滑到边界的区别仅作范围说明，不改成本书的紧支集稠密性结论。

\subsection{局部卷积}
\label{b1:subsec:220101}

设 \(\Omega\subseteq\mathbb R^d\) 为非空开集，\(k\in\mathbb N_0\)。仍使用\secref{b1:sec:2003}的非负偶磨光函数 \(\rho\)，其支集包含于 \(\overline B(0,1)\)，积分等于 \(1\)，并记 \(\rho_\varepsilon(x)=\varepsilon^{-d}\rho(x/\varepsilon)\)。对 \(u\in W^{k,p}_{\mathrm{loc}}(\Omega)\)，卷积 \(u*\rho_\varepsilon\) 只先定义在
\[
 \Omega_\varepsilon
 =\{\,x\in\Omega\mid\operatorname{dist}(x,\Omega^c)>\varepsilon\,\}.
\]
由\cref{b1:lem:sobolev-local-smoothing}，在任意固定内部相对紧开集上，这些卷积在有限 \(p\) 的 Sobolev 范数中趋于 \(u\)，并与不超过 \(k\) 阶的弱导数交换。不能在尚未定义 \(u\) 的地方替它指定零值，再假定这些导数关系仍成立。

若函数的非零部分已经留在内部紧集上，补零却没有这个问题。下面将“具有内部紧支集”明确解释为：存在 \(K\Subset\Omega\)，使函数在 \(\Omega\setminus K\) 上几乎处处为零。

\begin{lemma}\label{b1:lem:compact-sobolev-zero-extension}
设 \(1\leq p\leq\infty\)、\(u\in W^{k,p}(\Omega)\)，且 \(u\) 具有内部紧支集 \(K\)。在 \(\Omega\) 外将 \(u\) 补为零，所得函数记为 \(\widetilde u\)。则 \(\widetilde u\in W^{k,p}(\mathbb R^d)\)，并且
\begin{equation}\label{b1:eq:compact-sobolev-zero-extension}
 \partial^\alpha\widetilde u=\widetilde{\partial^\alpha u}
 \quad(|\alpha|\leq k),
 \quad
 \|\widetilde u\|_{W^{k,p}(\mathbb R^d)}
 =\|u\|_{W^{k,p}(\Omega)}.
\end{equation}
\end{lemma}
\begin{proof}
由弱导数的局部性及唯一性，\(\partial^\alpha u=0\) 几乎处处于 \(\Omega\setminus K\)。选取 \(\chi\in C_c^\infty(\Omega)\)，使其在 \(K\) 的一个邻域内等于 \(1\)；这样的截断由\cref{b1:lem:smooth-cutoff}给出。

任取 \(\varphi\in C_c^\infty(\mathbb R^d)\)。函数 \(\chi\varphi\) 是 \(\Omega\) 中的测试函数，而在 \(K\) 附近有 \(\partial^\alpha(\chi\varphi)=\partial^\alpha\varphi\)。因此
\[
 \begin{aligned}
 \int_{\mathbb R^d}\widetilde u\,\partial^\alpha\varphi\,\mathrm dx
 &=\int_\Omega u\partial^\alpha(\chi\varphi)\,\mathrm dx\\
 &=(-1)^{|\alpha|}\int_\Omega\partial^\alpha u\,\chi\varphi\,\mathrm dx\\
 &=(-1)^{|\alpha|}\int_{\mathbb R^d}
        \widetilde{\partial^\alpha u}\,\varphi\,\mathrm dx.
 \end{aligned}
\]
候选导数属于 \(L^p(\mathbb R^d)\)，于是得到弱导数等式。补零不改变任何分量的 \(L^p\) 范数，对有限 \(p\) 求 \(p\) 次和、对无穷端点取最大值，即得范数等式。
\end{proof}

特别地，若只有 \(u\in W^{k,p}_{\mathrm{loc}}(\Omega)\)，而 \(\psi\in C_c^\infty(\Omega)\)，则 \(\psi u\) 仍可使用这个引理。由\cref{b1:prop:weak-leibniz}，
\[
 \partial^\alpha(\psi u)
 =\sum_{\beta\leq\alpha}\binom\alpha\beta
       (\partial^{\alpha-\beta}\psi)(\partial^\beta u).
\]
每个光滑系数都在同一个内部紧集上有界，其支集外为零，所以每项全局属于 \(L^p(\Omega)\)。由此 \(\psi u\in W^{k,p}(\Omega)\)，且具有内部紧支集。这是随后使用单位分解的基本步骤。

\subsection{Friedrichs 磨光函数}
\label{b1:subsec:220102}

Friedrichs 磨光使用上述光滑核把函数换成邻域内的加权平均。\secref{b1:sec:2003}已经解释了这一操作为何使分布光滑；这里还要保留全部 Sobolev 范数估计。

\begin{proposition}\label{b1:prop:sobolev-friedrichs-smoothing}
若 \(u\in W^{k,p}(\mathbb R^d)\)、\(1\leq p\leq\infty\)，则
\(u_\varepsilon=\rho_\varepsilon*u\) 属于 \(C^\infty(\mathbb R^d)\cap W^{k,p}(\mathbb R^d)\)，且
\[
 \partial^\alpha u_\varepsilon
 =\rho_\varepsilon*\partial^\alpha u,
 \quad
 \|u_\varepsilon\|_{W^{k,p}(\mathbb R^d)}
 \leq\|u\|_{W^{k,p}(\mathbb R^d)}.
\]
若 \(p<\infty\)，还有 \(u_\varepsilon\to u\) 于 \(W^{k,p}(\mathbb R^d)\)。
\end{proposition}
\begin{proof}
光滑性和导数交换由\cref{b1:prop:distribution-convolution-smooth}成立。对 \(f\in L^p(\mathbb R^d)\)、\(p<\infty\)，在概率测度 \(\rho_\varepsilon(y)\,\mathrm dy\) 上使用 Hölder 不等式，得到
\[
 \bigl|\rho_\varepsilon*f(x)\bigr|^p
 \leq\int_{\mathbb R^d}\rho_\varepsilon(y)|f(x-y)|^p\,\mathrm dy.
\]
对 \(x\) 积分，Tonelli 定理与平移不变性给出 \(\|\rho_\varepsilon*f\|_p\leq\|f\|_p\)。当 \(p=\infty\) 时，非负核与单位积分直接给出同一估计。将它应用于有限多个导数，便得 Sobolev 范数界。

若 \(p<\infty\)，由\cref{b1:prop:approximate-identity-lp}，每个 \(\rho_\varepsilon*\partial^\alpha u\) 都在 \(L^p\) 中趋于 \(\partial^\alpha u\)。这些误差的 \(p\) 次方只有有限项，其和也趋零，得到所需范数收敛。
\end{proof}

无穷端点的结论只有范数不增，不能把最后一句也扩展到 \(p=\infty\)。\secref{b1:sec:2103}已用 \(|x|\) 的导数跳跃说明：即使在一个有界区间内部，光滑逼近也未必在 \(W^{1,\infty}\) 范数中成立。

\begin{proposition}\label{b1:prop:compact-sobolev-smooth-approximation}
设 \(1\leq p<\infty\)、\(u\in W^{k,p}(\Omega)\)，且 \(u\) 在 \(\Omega\setminus K\) 上几乎处处为零，其中 \(K\Subset\Omega\)。对任意包含 \(K\) 的开集 \(V\subseteq\Omega\)，当 \(\varepsilon\) 充分小时，
\[
 u_\varepsilon=\rho_\varepsilon*\widetilde u\in C_c^\infty(V),
 \quad
 \|u_\varepsilon-u\|_{W^{k,p}(\Omega)}\longrightarrow0.
\]
更精确地，\(\operatorname{supp}u_\varepsilon\subseteq K+\overline B(0,\varepsilon)\)，其中右侧表示两集合中向量的所有和。
\end{proposition}
\begin{proof}
由内部紧支集补零引理，\(\widetilde u\in W^{k,p}(\mathbb R^d)\)，故前一命题给出全空间范数收敛，限制到 \(\Omega\) 后仍收敛。若 \(\operatorname{dist}(x,K)>\varepsilon\)，卷积所读取的点都在 \(K\) 外，因此卷积为零。由其连续性，支集包含于紧集 \(K+\overline B(0,\varepsilon)\)。当 \(K\ne\varnothing\) 时，取
\(\varepsilon<\operatorname{dist}(K,V^c)\)，便使这个紧集包含于 \(V\)；空补集的距离仍约定为无穷。若 \(K=\varnothing\)，原函数及其卷积均为零，结论直接成立。
\end{proof}

这个命题既控制逼近误差，也允许预先指定一个支集必须留在其中的开集。仅有误差趋零还不足以拼接局部逼近：支集若在磨光后聚集到同一点，原来局部有限的函数族就可能失去局部有限性。

\subsection{Meyers–Serrin 定理}
\label{b1:subsec:220103}

\term[meyers-serrin-dingli]{Meyers--Serrin 定理}[Meyers--Serrin theorem]把内部磨光推广到任意开集上的全域范数逼近。它要求逼近函数在 \(\Omega\) 内光滑，不要求这些函数可以光滑延伸到边界外，也不要求它们具有紧支集。

\begin{theorem}[Meyers--Serrin]\label{b1:thm:meyers-serrin}
设 \(\Omega\subseteq\mathbb R^d\) 为非空开集，\(k\in\mathbb N_0\)、\(1\leq p<\infty\)。则 \(C^\infty(\Omega)\cap W^{k,p}(\Omega)\) 在实或复空间 \(W^{k,p}(\Omega)\) 中稠密。

更一般地，给定 \(u\in W^{k,p}_{\mathrm{loc}}(\Omega)\) 及 \(\eta>0\)，存在 \(v\in C^\infty(\Omega)\)，使 \(v-u\in W^{k,p}(\Omega)\) 且
\[
 \|v-u\|_{W^{k,p}(\Omega)}<\eta.
\]
后一陈述不要求 \(u\) 或 \(v\) 各自属于全局 \(W^{k,p}(\Omega)\)。
\end{theorem}
\begin{proof}
证明较一般的陈述。沿用\cref{b1:thm:smooth-partition-unity}证明中的紧集穷竭，记
\[
 K_j=\{\,x\in\Omega\mid |x|\leq j,
                 \operatorname{dist}(x,\Omega^c)\geq1/j\,\},
 \quad K_0=K_{-1}=\varnothing,
\]
并置 \(V_j=\operatorname{int}K_{j+1}\setminus K_{j-2}\)。紧层 \(K_j\setminus\operatorname{int}K_{j-1}\) 覆盖 \(\Omega\)，且各自包含于 \(V_j\)，故这些开集也覆盖 \(\Omega\)。它们局部有限：任意一点有一个包含在某个 \(\operatorname{int}K_N\) 中的邻域，所有 \(j\geq N+2\) 的 \(V_j\) 都避开它。

该单位分解证明在每一层构造有限多个函数，其支集紧含于 \(V_j\)。把这些函数排成 \((\psi_i)_{i\geq1}\)，略去零函数，记第 \(i\) 个函数所属层为 \(j(i)\)。于是
\[
 \psi_i\geq0,
 \quad\sum_i\psi_i=1,
 \quad\operatorname{supp}\psi_i\Subset V_{j(i)},
\]
且每层只出现有限次。令 \(w_i=\psi_i u\)，并在 \(\Omega\) 外补零。前面的乘积讨论及补零引理保证 \(\widetilde w_i\in W^{k,p}(\mathbb R^d)\)。

对每个 \(i\)，使用\cref{b1:prop:compact-sobolev-smooth-approximation}，选取半径 \(\varepsilon_i>0\)，使
\[
 v_i=\rho_{\varepsilon_i}*\widetilde w_i\in C_c^\infty\bigl(V_{j(i)}\bigr),
 \quad
 \|v_i-w_i\|_{W^{k,p}(\Omega)}<\eta2^{-i-1}.
\]
这里先要求 \(\varepsilon_i<\operatorname{dist}\bigl(\operatorname{supp}\psi_i,V_{j(i)}^c\bigr)\)，再进一步缩小半径以达到误差条件。两个要求可以同时实现。

函数族 \((v_i)\) 仍局部有限。事实上，每个紧集只遇到有限层 \(V_j\)，而每层只有有限项。因此
\[
 v=\sum_{i=1}^\infty v_i
\]
在每个点的某个邻域内都是有限个光滑函数的和，故 \(v\in C^\infty(\Omega)\)。

剩下的全域估计应对误差求和。记 \(e_i=v_i-w_i\)。因为
\[
 \sum_i\|e_i\|_{W^{k,p}(\Omega)}\leq\eta/2,
\]
由\cref{b1:thm:sobolev-complete}及\ref{b1:prop:banach-series-criterion}，级数 \(\sum_i e_i\) 在 \(W^{k,p}(\Omega)\) 中收敛到某个 \(e\)，且 \(\|e\|_{W^{k,p}}\leq\eta/2\)。在任意内部相对紧开集上，\(v_i\) 和 \(w_i\) 都只有有限个非零项；足够长的部分和因而几乎处处等于
\[
 \sum_i(v_i-w_i)=v-u\sum_i\psi_i=v-u.
\]
另一方面，部分和在全域 \(L^p\) 中趋于 \(e\)，从而也在这个开集上趋于 \(e\)。故 \(e=v-u\) 几乎处处于该开集。再取可数个这样的开集穷竭 \(\Omega\)，得到全域等式与所需误差界。

若原来的 \(u\in W^{k,p}(\Omega)\)，则 \(v=u+e\) 也属于该空间。对每个正整数 \(n\) 取 \(\eta=1/n\)，便得到定理中的稠密性。
\end{proof}

证明没有宣称 \(\sum_{i=1}^N\psi_i u\) 在全域 Sobolev 范数中趋于 \(u\)。单位分解的导数可能随位置增大；局部有限只保证局部运算合法，并不自动给出全域范数收敛。逐项磨光后的误差另有可和估计，完备性用在这组误差上。

\subsection{\texorpdfstring{\(C_c^\infty\)}{Cc∞} 的稠密性问题}
\label{b1:subsec:220104}

Meyers--Serrin 定理产生的光滑函数可能一直延伸到边界附近或无穷远处。紧支集是额外要求，需要单独判断。

\begin{theorem}\label{b1:thm:whole-space-sobolev-density}
对 \(k\in\mathbb N_0\)、\(1\leq p<\infty\)，空间 \(C_c^\infty(\mathbb R^d)\) 在 \(W^{k,p}(\mathbb R^d)\) 中稠密。
\end{theorem}
\begin{proof}
由光滑截断引理，取 \(0\leq\chi\leq1\)，使 \(\chi=1\) 于 \(\overline B(0,1)\) 的一个邻域，且 \(\chi\in C_c^\infty\bigl(B(0,2)\bigr)\)。对 \(R\geq1\)，置 \(\chi_R(x)=\chi(x/R)\)。任取 \(u\in W^{k,p}(\mathbb R^d)\)。由 Leibniz 公式，对 \(|\alpha|\leq k\)，
\[
 \begin{aligned}
 \partial^\alpha(\chi_Ru-u)
 &=(\chi_R-1)\partial^\alpha u\\
 &\quad+\sum_{\substack{\beta\leq\alpha\\|\beta|>0}}
   \binom\alpha\beta R^{-|\beta|}
   (\partial^\beta\chi)(x/R)\partial^{\alpha-\beta}u.
 \end{aligned}
\]
第一项的 \(L^p\) 范数不超过
\(\|\mathbf1_{\mathbb R^d\setminus B(0,R)}\partial^\alpha u\|_p\)，由有限 \(p\) 的可积尾部趋于零。余下每项的范数不超过
\[
 \binom\alpha\beta R^{-|\beta|}
 \|\partial^\beta\chi\|_\infty\cdot\|\partial^{\alpha-\beta}u\|_p,
\]
也趋于零。有限个多重指标合起来给出 \(\chi_Ru\to u\) 于 \(W^{k,p}(\mathbb R^d)\)。

给定 \(\eta>0\)，先取 \(R\) 使截断误差小于 \(\eta/2\)。函数 \(\chi_Ru\) 具有紧支集，再由\cref{b1:prop:compact-sobolev-smooth-approximation}取 \(v\in C_c^\infty(\mathbb R^d)\)，使 \(\|v-\chi_Ru\|_{W^{k,p}}<\eta/2\)。三角不等式完成逼近。
\end{proof}

这里的截断只切掉无穷远处的尾部，并且 \(\chi_R\) 的非零阶导数带有衰减因子 \(R^{-|\beta|}\)。若试图在一般域的边界附近作越来越陡的截断，其导数反而可能增大，不能原样复制这个估计。

例如在 \(I=(0,1)\) 上取 \(u=1\)。它属于所有有限 \(p\) 的 \(W^{1,p}(I)\)。对任意 \(\varphi\in C_c^\infty(I)\)，函数 \(1-\varphi\) 的连续代表在两个端点都等于 \(1\)，因而由\eqref{b1:eq:interval-sobolev-sup}，
\[
 \begin{aligned}
 1&\leq\|1-\varphi\|_{L^1(I)}+\|\varphi'\|_{L^1(I)}\\
  &\leq\|1-\varphi\|_{L^p(I)}+\|\varphi'\|_{L^p(I)}\\
  &\leq2^{1-1/p}\|1-\varphi\|_{W^{1,p}(I)}.
 \end{aligned}
\]
最后一步是两项有限序列的 Hölder 不等式，也包括 \(p=1\)。因此这些紧支集函数到常数 \(1\) 的距离至少为 \(2^{1/p-1}>0\)。同一估计还排除了任意 \(k\geq1\) 时在 \(W^{k,p}(I)\) 中的逼近，因为高阶范数也控制其中的零阶和一阶项。问题来自边界附近强制为零的要求，不是域内光滑性不足；常数 \(1\) 本身已经光滑。

零阶情形应另作区分。对任意开集 \(\Omega\) 和有限 \(p\)，\(C_c^\infty(\Omega)\) 在 \(L^p(\Omega)\) 中仍稠密。事实上，取内部紧集穷竭 \(K_j\) 及 \(0\leq\chi_j\leq1\)、\(\chi_j\in C_c^\infty(\Omega)\)，使 \(\chi_j=1\) 于 \(K_j\)。则 \(\chi_j u\to u\) 于 \(L^p\)，这是控制收敛定理的直接应用；再对每个 \(\chi_j u\) 使用 \(k=0\) 的内部紧支集磨光命题即可。只控制函数值时，边界层的 \(L^p\) 误差可以消失；同时控制导数后，这个边界层便可能留下不可忽略的代价。下一节将把能够由紧支集光滑函数逼近的 Sobolev 函数单独组成一个闭子空间。

% !TeX root = ../../Book1.tex
\section{\texorpdfstring{\(W_0^{1,p}\)}{W₀(1,p)}}
\label{b1:sec:2202}

域内光滑逼近不要求近似函数在边界附近消失。若近似函数还须有紧支集，上一节的一维反例已经表明，得到的闭包可能是真子空间。本节就以这个闭包定义零边界空间，再说明它与补零、截断以及经典端点值之间的关系。对一般开集，定义本身不预设迹算子存在。

\subsection{定义}
\label{b1:subsec:220201}

\begin{definition}\label{b1:def:sobolev-zero-boundary}
设 \(\Omega\subset\mathbb R^d\) 为非空开集、\(1\le p\le\infty\)。定义
\[
 W_0^{1,p}(\Omega)
 =\overline{C_c^\infty(\Omega)}^{\,W^{1,p}(\Omega)}.
\]
也就是说，\(u\in W_0^{1,p}(\Omega)\) 当且仅当存在 \(\varphi_j\in C_c^\infty(\Omega)\)，使 \(\|u-\varphi_j\|_{W^{1,p}(\Omega)}\to0\)。另记 \(H_0^1(\Omega)=W_0^{1,2}(\Omega)\)。
\end{definition}
\symidx{\(W_0^{1,p}(\Omega)\)：紧支集光滑函数在\(W^{1,p}\)中的闭包}
\symidx{\(H_0^1(\Omega)\)：\(W_0^{1,2}(\Omega)\)}

它是 \(W^{1,p}\) 的闭线性子空间，所以关于原范数完备；\(H_0^1\) 关于原 \(H^1\) 内积为 Hilbert 空间。若 \(1<p<\infty\)，闭子空间的自反性又给出 \(W_0^{1,p}\) 自反。这些结论都来自定义中的闭包，不需要对边界作几何假设。

下标 \(0\) 通常用来表达零边界条件。不过在这个定义中，“零”首先落实为可用内部紧支集函数逼近，而不是给每个边界点指定一个数。一般 Sobolev 元素仅是几乎处处等价类，边界点甚至不属于定义域；若没有额外的代表定理，写下逐点边界值就还没有确定一个运算。

\(p=\infty\) 时，本书也严格使用上述范数闭包。它通常比“有零边界值的 Lipschitz 函数”更小，不能把有限 \(p\) 的稠密性结论直接搬到这个端点。需要时，应明确其他文献所用 \(W_0^{1,\infty}\) 是否采用同一定义。

\subsection{零边界条件}
\label{b1:subsec:220202}

先把不涉及边界点取值、却总是有效的补零结论证明出来。对 \(u\) 在 \(\Omega\) 上的可测代表，令
\[
 E_0u(x)=
 \begin{cases}
 u(x),&x\in\Omega,\\
 0,&x\notin\Omega.
 \end{cases}
\]
这个操作在 \(L^p\) 等价类上始终有定义并保持范数，困难只在于补零后的导数。

\begin{proposition}\label{b1:prop:w0-zero-extension}
若 \(u\in W_0^{1,p}(\Omega)\)，则 \(E_0u\in W^{1,p}(\mathbb R^d)\)，并且
\[
 \partial_i(E_0u)=E_0(\partial_i u),\qquad
 \|E_0u\|_{W^{1,p}(\mathbb R^d)}=\|u\|_{W^{1,p}(\Omega)}.
\]
因此 \(E_0\) 是一个线性等距嵌入，结论包括 \(p=\infty\)。
\end{proposition}
\begin{proof}
取定义中的 \(\varphi_j\)。每个 \(E_0\varphi_j\) 都属于 \(C_c^\infty(\mathbb R^d)\)，而且经典导数在域内外直接给出
\[
 \partial_i(E_0\varphi_j)=E_0(\partial_i\varphi_j).
\]
因此 \(\|E_0\varphi_j-E_0\varphi_k\|_{W^{1,p}(\mathbb R^d)}
=\|\varphi_j-\varphi_k\|_{W^{1,p}(\Omega)}\)。由完备性，它们在全空间的 \(W^{1,p}\) 中收敛。另一方面，\(L^p\) 补零等距性表明其零阶极限为 \(E_0u\)，各一阶分量的极限为 \(E_0(\partial_i u)\)。这就同时证明了导数公式及范数等式。
\end{proof}

反向在任意开集上不成立。令
\[
 \Omega=(-1,0)\cup(0,1),\qquad u(x)=1-|x|\quad(x\in\Omega).
\]
补零后所得等价类由全实轴上的帐篷函数 \((1-|x|)^+\) 表示；仅在原点处填的值不同，不影响等价类，所以 \(E_0u\in W^{1,p}(\mathbb R)\) 对全部 \(p\) 成立。然而若 \(\varphi_j\in C_c^\infty(\Omega)\) 依 \(W^{1,p}(\Omega)\) 逼近 \(u\)，限制到 \((0,1)\)，由\secref{b1:sec:2104}的一维上确界估计，绝对连续代表必须在 \([0,1]\) 上一致收敛。近似函数在 \(0\) 处都为零，极限代表在 \(0\) 处却为一，矛盾。

这里补零只改变一个内部缺失点，积分看不见这个点；但沿两侧区间逼近时，端点行为仍受到导数范数控制。因此“补零没有产生奇异导数”与“能够在域内由紧支集函数逼近”是两条不同的条件。

无穷端点还有另一种障碍。在 \((0,1)\) 上，\(u(x)=x(1-x)\) 光滑且两个端点值都为零，但它不属于 \(W_0^{1,\infty}(0,1)\)。事实上，每个 \(\varphi\in C_c^\infty(0,1)\) 在零点附近恒为零，故 \(\varphi'=0\) 在该邻域；而 \(u'(x)=1-2x\) 趋于一。因此
\[
 \|u'-\varphi'\|_{L^\infty(0,1)}\ge1.
\]
这不是端点函数值的问题，而是强无穷范数还要求导数在靠近端点时也能被一致控制。

\subsection{截断逼近}
\label{b1:subsec:220203}

以后在积分恒等式中使用 Sobolev 函数作测试时，常需先控制它的幅值。有限 \(p\) 下可以在零边界空间内完成这一步。

\begin{proposition}\label{b1:prop:w0-truncation}
设 \(1\le p<\infty\)，实值 \(u\in W_0^{1,p}(\Omega)\)。则对每个 \(M>0\)，截断 \(T_M(u)\) 仍属于 \(W_0^{1,p}(\Omega)\)，且
\[
 T_M(u)\longrightarrow u\quad\text{于 }W^{1,p}(\Omega)
 \quad(M\to\infty).
\]
正部、负部和绝对值也保持 \(W_0^{1,p}\)。若另外 \(|u|\le M\) 几乎处处，则可以选取 \(\psi_j\in C_c^\infty(\Omega)\)，使 \(|\psi_j|\le M\) 且 \(\psi_j\to u\) 于 \(W^{1,p}\)。
\end{proposition}
\begin{proof}
取 \(\varphi_j\in C_c^\infty(\Omega)\) 逼近 \(u\)。固定 \(M\)，先说明
\[
 T_M(\varphi_j)\longrightarrow T_M(u)\quad\text{于 }W^{1,p}(\Omega).
\]
零阶收敛由 \(T_M\) 的 \(1\)-Lipschitz 性得到。一阶导数使用上一章的截断公式，分解为
\[
 \begin{aligned}
 \partial_iT_M(\varphi_j)-\partial_iT_M(u)
 &=\mathbf1_{\{|\varphi_j|<M\}}(\partial_i\varphi_j-\partial_i u)\\
 &\quad+
  \bigl(\mathbf1_{\{|\varphi_j|<M\}}-\mathbf1_{\{|u|<M\}}\bigr)\partial_i u.
 \end{aligned}
\]
第一项在 \(L^p\) 中趋零。沿任意一条 \(\varphi_j\to u\) 几乎处处的子列，第二项在 \(|u|\ne M\) 处几乎处处趋零；在 \(|u|=M\) 上，水平集梯度为零。控制收敛于是给出该项的 \(L^p\) 收敛。对任意子列再抽取这样的子列，可把结论传回全列，正如习题\ref{b1:ex:21-positive-part-continuity}中的论证。

每个 \(T_M(\varphi_j)\) 都是内部紧支集的 \(W^{1,p}\) 函数。由命题\ref{b1:prop:compact-sobolev-smooth-approximation}，它属于 \(W_0^{1,p}\)。再用闭性，得到 \(T_M(u)\in W_0^{1,p}\)。当 \(M\to\infty\) 时的强收敛已由\eqref{b1:eq:sobolev-truncation-convergence}证明。正部、负部和绝对值用相同的紧支集逼近，并使用上一章相应的强连续性，即可得到。

若 \(|u|\le M\)，则 \(T_M(u)=u\)。对每个 \(T_M(\varphi_j)\) 选择足够小的非负单位积分卷积，使支集仍留在 \(\Omega\)，且 \(W^{1,p}\) 误差小于 \(1/j\)。卷积不增加幅值上界，故得到所需的 \(\psi_j\)。
\end{proof}

最后一段同时说明了一个常用事实：内部紧支集的 \(W^{1,p}\) 函数在有限 \(p\) 下属于 \(W_0^{1,p}\)。无穷端点不能以同一理由推出结论，因为磨光并不保证 \(W^{1,\infty}\) 范数收敛；一个内部尖点已经可能构成障碍。这里对截断保持性的断言也只在有限 \(p\) 范围内使用。

\subsection{与迹的关系}
\label{b1:subsec:220204}

在有界区间 \(I=(a,b)\) 上，上一章已经给每个 \(W^{1,p}(I)\) 元素选出了唯一的闭区间绝对连续代表。于是端点取值
\[
 \gamma_Iu=(\widetilde u(a),\widetilde u(b))
\]
是良定义的线性映射；由\eqref{b1:eq:interval-sobolev-sup}及 Hölder 不等式，它对 \(W^{1,p}\) 范数连续。这是一般迹算子的一维原型。

\begin{theorem}\label{b1:thm:interval-zero-boundary}
若 \(I=(a,b)\) 有界、\(1\le p<\infty\)，则
\[
 W_0^{1,p}(I)
 =\{\,u\in W^{1,p}(I):\widetilde u(a)=\widetilde u(b)=0\,\}
 =\ker\gamma_I.
\]
\end{theorem}
\begin{proof}
一个方向由 \(\gamma_I\) 连续且在 \(C_c^\infty(I)\) 上为零立即得到。反过来，设两个端点值均为零。对充分小的 \(\varepsilon>0\)，选 \(0\le\eta_\varepsilon\le1\)，使其在距端点不超过 \(\varepsilon\) 处为零、在 \([a+2\varepsilon,b-2\varepsilon]\) 上为一，且
\[
 \eta_\varepsilon\in C_c^\infty(I),\qquad
 |\eta_\varepsilon'|\le C/\varepsilon.
\]
乘积 \(u_\varepsilon=\eta_\varepsilon u\) 为内部紧支集的 Sobolev 函数，因此属于 \(W_0^{1,p}(I)\)。函数差 \((1-\eta_\varepsilon)u\) 与导数项 \((1-\eta_\varepsilon)u'\) 的 \(L^p\) 范数都因边界层缩小而趋零。需要单独估计的是 \(\eta_\varepsilon'u\)。

由 \(u(a)=0\) 和绝对连续性，对 \(a<x<a+2\varepsilon\)，Hölder 不等式给出
\[
 |u(x)|^p
 \le (x-a)^{p-1}\int_a^x|u'(t)|^p\,\mathrm dt.
\]
\(p=1\) 时这就是积分三角不等式。再次积分并用 Tonelli 定理，得到
\[
 \varepsilon^{-p}\int_a^{a+2\varepsilon}|u(x)|^p\,\mathrm dx
 \le 2^p\int_a^{a+2\varepsilon}|u'(t)|^p\,\mathrm dt
 \longrightarrow0.
\]
在右端点使用 \(u(b)=0\) 得到同样估计。因此 \(\eta_\varepsilon'u\to0\) 于 \(L^p\)，所以 \(u_\varepsilon\to u\) 于 \(W^{1,p}(I)\)。利用 \(W_0^{1,p}\) 的闭性完成证明。
\end{proof}

这个证明显示了零端点值的具体作用：它把 \(u\) 在薄边界层中的大小变成 \(u'\) 在同一层上的积分，从而抵消截断导数中的 \(\varepsilon^{-1}\)。仅仅知道 \(u\) 有界，不足以完成这一步。前面的 \(x(1-x)\) 例子则说明，定理的有限 \(p\) 范围不可直接放宽到无穷端点。

对有界开集，还有一个不要求边界规则的充分条件。若实值 \(u\in W^{1,p}(\Omega)\cap C(\overline\Omega)\)、\(p<\infty\)，并且 \(u=0\) 在 \(\partial\Omega\) 上，则 \(u\in W_0^{1,p}(\Omega)\)。为此取
\[
 v_\varepsilon=u-T_\varepsilon(u).
\]
由连续性与零边界值，\(v_\varepsilon\) 的支集包含于紧集 \(\{x\in\overline\Omega:|u(x)|\ge\varepsilon\}\Subset\Omega\)，所以它属于 \(W_0^{1,p}\)。而 \(T_\varepsilon(u)\to0\) 于 \(L^p\)，其弱梯度为 \(\mathbf1_{\{|u|<\varepsilon\}}\nabla u\)，由水平集梯度为零及控制收敛也趋于零。故 \(v_\varepsilon\to u\) 于 \(W^{1,p}\)。复值情形对实部、虚部分别处理即可。

一般 Sobolev 元素未必具有这样的连续代表；一般边界也未必允许只用逐点限制描述其行为。下一章将对适当的区域构造有界迹算子，并在明确条件下辨认其核。在此之前，\(W_0^{1,p}\) 始终以紧支集光滑闭包为定义，而不是以尚未定义的边界取值替换它。

% !TeX root = ../../Book1.tex
\section{延拓算子}
\label{b1:sec:2203}

在开集内部磨光时，可以把卷积半径选得足够小；接近边界以后，这种选择不再给出同一个全域函数。延拓提供了另一条路线：先在开集外补出函数，同时控制弱导数，再在整个空间中使用卷积。本节只处理一阶空间，且保留两个端点 \(p=1\) 与 \(p=\infty\)。

有限图表、反射和单位分解的几何组织，可对照 Di Nezza、Palatucci 与 Valdinoci 的《Hitchhiker's guide to the fractional Sobolev spaces》\cite[Section 5]{DiNezza2012Hitchhikers}。该文讨论分数阶空间；下面的一阶证明则使用前章的沿线绝对连续代表与双 Lipschitz 换元，不由分数阶结论直接推出。
% 实读来源：DiNezza2012Hitchhikers，作者公开副本 arXiv:1104.4345，
% 第5节 Lemmas 5.1--5.3 与 Theorem 5.4 的完整证明。
% 原文范围为 0<s<1、1<=p<infty；这里只借鉴图表、截断和反射的组织，
% 一阶及 p=infty 的论证由本书21.3--21.4接口独立完成。

\subsection{半空间反射}
\label{b1:subsec:220301}

记上半空间为 \(H_+=\mathbb R^{d-1}\times(0,\infty)\)，点写成 \(x=(x',t)\)。反射时保留函数值，法向导数则改变符号。关键不是这个形式计算，而是证明两个半空间的弱导数能跨过中间的超平面拼在一起。

\begin{theorem}[半空间偶反射]\label{b1:thm:half-space-sobolev-extension}
设 \(1\le p\le\infty\)。对 \(u\in W^{1,p}(H_+)\)，在 \(t\ne0\) 处规定
\[
 (E_+u)(x',t)=u(x',|t|).
\]
该公式定义一个从 \(W^{1,p}(H_+)\) 到 \(W^{1,p}(\mathbb R^d)\) 的线性映射，其弱导数为
\begin{equation}\label{b1:eq:half-space-reflection-derivatives}
 \partial_i(E_+u)(x',t)=
 \begin{cases}
  (\partial_i u)(x',|t|),&1\le i<d,\\
  \operatorname{sgn}(t)(\partial_d u)(x',|t|),&i=d.
 \end{cases}
\end{equation}
此外，按本书的 Sobolev 范数，
\[
 \|E_+u\|_{W^{1,p}(\mathbb R^d)}
 =2^{1/p}\|u\|_{W^{1,p}(H_+)}\quad(1\le p<\infty),
\]
而 \(p=\infty\) 时两侧范数相等。
\end{theorem}
\begin{proof}
先取\cref{b1:thm:sobolev-acl}给出的共同 Borel 代表。由 Fubini 定理及有界盒上的 Hölder 不等式，对几乎每个 \(x'\)，函数 \(u(x',\cdot)\) 与 \(\partial_d u(x',\cdot)\) 在每个 \((0,R)\) 上均可积。只需先取正整数 \(R\)，便可使这里的例外集与 \(R\) 无关。沿线的绝对连续性与\cref{b1:prop:sobolev-interval-representative}于是给出端点值 \(a(x')\)，并有
\[
 u(x',t)=a(x')+\int_0^t\partial_d u(x',r)\,\mathrm dr
 \quad(0\le t\le R).
\]
两侧反射使用同一个 \(a(x')\)，故所得函数在 \([-R,R]\) 上绝对连续，导数为\eqref{b1:eq:half-space-reflection-derivatives}的法向分量。端点值可取为 \(u(x',1/k)\) 的极限；在有限极限不存在处取零，从而仍得到 Borel 代表。

对其他坐标方向，\(t\ne0\) 的直线只是原来水平直线的反射，其绝对连续性及切向导数公式均得保留。超平面 \(t=0\) 所在的整层直线，在这些方向的横向参数中为零测集，可以一并排除。因此新函数有同一个适用于全部方向的沿线绝对连续代表。

反射保持 Lebesgue 测度，故新函数及候选弱导数都在 \(L^p(\mathbb R^d)\) 中。由 ACL 的反向判据，它们确实满足弱导数关系。对有限 \(p\)，函数和每个一阶导数的 \(p\) 次积分都恰好加倍；对 \(p=\infty\)，本性上确界不变，这就给出范数等式。最后，反射把零测集映为零测集，公式不依赖 \(u\) 的代表；超平面上的赋值不改变等价类。因而算子在等价类层面是线性的。
\end{proof}

如果把 \(u(t)=e^{-t}\) 从 \((0,\infty)\) 直接补零到负半轴，原点两侧的值不吻合，分布导数便含有点质量。偶反射得到 \(e^{-|t|}\)，两侧连续接合，其一阶弱导数只有普通函数部分。这里需要的是函数值接合，而不是要求法向导数也在边界两侧相等。

\subsection{Lipschitz 区域}
\label{b1:subsec:220302}

半空间模型可以在边界附近经过坐标变换使用。为此，边界不必具有连续法向，但须能在足够小的坐标柱中写成一张 Lipschitz 图像。

\begin{definition}\label{b1:def:lipschitz-domain}
设 \(\Omega\subset\mathbb R^d\) 为开集。若每个 \(a\in\partial\Omega\) 都有刚性坐标，使 \(a\) 的坐标为零，并有 \(R,H>0\) 及 Lipschitz 函数
\(\varphi:(-R,R)^{d-1}\rightarrow(-H,H)\)，满足 \(\varphi(0)=0\) 和
\[
 \Omega\cap C=\{\,(z,s)\in C\mid s>\varphi(z)\,\},
 \quad C=(-R,R)^{d-1}\times(-H,H),
\]
则称 \(\Omega\) 为\term[lipschitz-quyu]{Lipschitz 区域}[Lipschitz domain]。本章不要求它连通；需要有界性时另行说明。维数为一时，横向坐标省略，条件就是边界附近为区间的一端。
\end{definition}

例如平面第一象限在角点旋转坐标后可写成 \(s>|z|\)，因此角点本身不妨碍这一条件。有限个互不接触的闭球的内部之并，也符合本章不要求连通的约定。这里允许折角，但不允许在同一局部图柱中把开集放在图像的两侧。

下面固定图表尺寸，以便在反射以后仍有空间补零。设图函数的 Lipschitz 常数为 \(L\)，选取 \(\rho>0\)，使
\[
 2\rho<R,\quad 2L\sqrt{d-1}\rho<H/4,
 \quad h=H/4.
\]
若 \(L=0\) 或 \(d=1\)，第二个限制自动满足。取对称盒
\[
 \begin{split}
 D&=(-2\rho,2\rho)^{d-1}\times(-2h,2h),\\
 D_1&=(-\rho,\rho)^{d-1}\times(-h,h),
 \end{split}
\]
并在上述刚性坐标中规定
\[
 F(z,t)=(z,\varphi(z)+t),\quad
 F^{-1}(z,s)=(z,s-\varphi(z)).
\]
由于 \(|\varphi(z)|<H/4\) 且 \(|t|<H/2\)，整个 \(F(D)\) 都留在原图柱内。映射 \(F\) 与 \(F^{-1}\) 的 Lipschitz 常数至多为 \(1+L\)。回到原坐标以后，刚性变换不破坏这一控制。

记 \(U=F(D)\)、\(V=F(D_1)\) 及 \(D^+=D\cap H_+\)。于是 \(\overline V\subset U\)，且
\[
 \Omega\cap U=F(D^+).
\]
两个盒都关于 \(t=0\) 对称；内盒中的集合反射后仍在内盒中。后面将截断函数放在 \(V\) 内，却在较大的 \(U\) 中完成延拓，内外两层之间的余量用来避免侧面和顶面上的跳跃。

\subsection{延拓算子}
\label{b1:subsec:220303}

\begin{definition}\label{b1:def:sobolev-extension-operator}
给定开集 \(\Omega\) 及 \(1\le p\le\infty\)，若线性映射
\[
 E:W^{1,p}(\Omega)\longrightarrow W^{1,p}(\mathbb R^d)
\]
满足 \((Eu)|_\Omega=u\) 几乎处处，则称它为该空间的\term[yantuo-suanzi]{延拓算子}[extension operator]。若还有与 \(u\) 无关的常数 \(C\)，使
\[
 \|Eu\|_{W^{1,p}(\mathbb R^d)}
 \le C\|u\|_{W^{1,p}(\Omega)},
\]
则称这个延拓算子有界。
\end{definition}

先处理落在一个边界图表内的部分。此时截断函数在物理空间中光滑；拉回后只保证 Lipschitz，所以先作光滑乘法、再作坐标变换，可以直接使用已经建立的两条规则。

\begin{lemma}\label{b1:lem:localized-graph-extension}
固定上一小节满足 \(\overline V\Subset U\) 的内外图表及 \(\eta\in C_c^\infty(V)\)。存在与 \(p\) 无关的线性构造 \(E_\eta\)，使每个 \(u\in W^{1,p}(\Omega)\) 都满足
\[
 E_\eta u\in W^{1,p}(\mathbb R^d),\quad
 (E_\eta u)|_\Omega=\eta u,
\]
并有
\[
 \|E_\eta u\|_{W^{1,p}(\mathbb R^d)}
 \le C_{\eta,F,p}\|u\|_{W^{1,p}(\Omega\cap U)}.
\]
延拓的支集包含在一个只依赖于 \(\eta,F\) 的固定紧集内，该紧集包含于 \(U\)。
\end{lemma}
\begin{proof}
光滑乘法的弱 Leibniz 公式给出
\[
 \partial_i(\eta u)=\eta\partial_i u+(\partial_i\eta)u.
\]
因此乘法在这里的 Sobolev 范数下有界，不要求 \(u\) 本性有界。由\cref{b1:thm:sobolev-bilipschitz-change}，函数
\[
 v=(\eta u)\circ F\in W^{1,p}(D^+)
\]
也有相应范数控制。它在 \(D^+\setminus\overline{D_1}\) 上为零，其弱导数在这一开集上同样为零。

将 \(v\) 在 \(H_+\setminus D^+\) 上补零，记为 \(\widehat v\)。这里的支集可能贴着 \(t=0\)，不能把它当作紧含于 \(D^+\) 的内部函数。不过，开集
\[
 D^+,\quad H_+\setminus\overline{D_1}
\]
覆盖整个 \(H_+\)，且在交集上，\(v\) 与零函数及其弱导数都相同。任取 \(H_+\) 上的测试函数，利用\cref{b1:thm:smooth-partition-unity}将它分解为有限个测试函数之和，每个分量支集落在上述某个开集中。在各分量上分别使用 \(v\) 或零函数的弱导数恒等式，再求和，便得到 \(\widehat v\) 的弱导数。候选导数正是原导数的零延拓，因此
\[
 \|\widehat v\|_{W^{1,p}(H_+)}=\|v\|_{W^{1,p}(D^+)}.
\]

现在作半空间偶反射 \(w=E_+\widehat v\)。令
\[
 S=F^{-1}(\operatorname{supp}\eta)\cap\{\,t\ge0\,\},
 \quad \mathcal R(z,t)=(z,-t).
\]
集合 \(S\) 是 \(D_1\) 中的紧集，\(w\) 在 \(S\cup\mathcal R(S)\) 之外几乎处处为零。因此反射后的支集仍紧含于 \(D_1\)。再在 \(U\) 上把 \(w\) 推回物理空间：
\[
 e=w\circ F^{-1}.
\]
双 Lipschitz 换元保证 \(e\in W^{1,p}(U)\)，且它的支集包含于固定紧集
\(F\bigl(S\cup\mathcal R(S)\bigr)\Subset U\)。此时可用\cref{b1:lem:compact-sobolev-zero-extension}把 \(e\) 补零到整个空间，所得函数就是 \(E_\eta u\)。

在 \(\Omega\cap U\) 上，拉回坐标的最后一分量为正，所以反射并未改变原函数，故 \(E_\eta u=\eta u\)；在 \(\Omega\setminus U\) 上两者都为零。各次乘法、复合、反射与补零都是固定的线性操作。依次使用它们的范数界，就得到引理中的估计。
\end{proof}

\subsection{有界延拓}
\label{b1:subsec:220304}

有界性使边界能够由有限个内图表覆盖，从而把局部延拓合成一个全域算子。

\begin{theorem}[Lipschitz 区域的一阶延拓]\label{b1:thm:lipschitz-domain-extension}
设 \(\Omega\subset\mathbb R^d\) 为有界 Lipschitz 区域。可以固定同一套线性延拓公式，使它对每个 \(1\le p\le\infty\) 都定义有界算子
\[
 E:W^{1,p}(\Omega)\longrightarrow W^{1,p}(\mathbb R^d),
 \quad (Eu)|_\Omega=u,
\]
并满足
\begin{equation}\label{b1:eq:lipschitz-domain-extension-bound}
 \|Eu\|_{W^{1,p}(\mathbb R^d)}
 \le C_{\Omega,p}\|u\|_{W^{1,p}(\Omega)}.
\end{equation}
所有 \(Eu\) 均可取支集包含于同一个固定紧集。
\end{theorem}
\begin{proof}
由边界的紧性，选有限个内图表 \(V_1,\ldots,V_N\) 覆盖 \(\partial\Omega\)，对应外图表记为 \(U_j\)。集合
\[
 K_0=\overline\Omega\setminus\bigcup_{j=1}^N V_j
\]
是 \(\Omega\) 内的紧集。选开集 \(V_0\) 包含 \(K_0\)，并使 \(\overline{V_0}\Subset\Omega\)。若 \(K_0\) 为空，可省去这一内部项。

在开集 \(N_0=\bigcup_{j=0}^N V_j\) 上取从属于这有限覆盖的光滑单位分解，再取 \(\chi\in C_c^\infty(N_0)\)，使它在 \(\overline\Omega\) 的一个邻域 \(N_1\) 中等于 \(1\)。单位分解局部有限，所以只有有限个分解函数的支集与 \(\operatorname{supp}\chi\) 相交。将这些函数乘以 \(\chi\)，并按所属的 \(V_j\) 合并，得到
\[
 \eta_j\in C_c^\infty(V_j),\quad
 \sum_{j=0}^N\eta_j=1
 \quad(x\in N_1).
\]
这一步只依赖区域和图表，不依赖待延拓的函数。

内部函数 \(\eta_0u\) 由\cref{b1:lem:compact-sobolev-zero-extension}补零，记为 \(e_0\)。对每个边界项应用\cref{b1:lem:localized-graph-extension}，记其延拓为 \(E_{\eta_j}u\)，并规定
\[
 Eu=e_0+\sum_{j=1}^N E_{\eta_j}u.
\]
限制到 \(\Omega\) 后，各项之和为 \(\sum_j\eta_j u=u\)。所有构造均线性，所以 \(E\) 线性。用三角不等式及各局部估计，得到
\[
 \begin{split}
 \|Eu\|_{W^{1,p}(\mathbb R^d)}
 &\le\|e_0\|_{W^{1,p}(\mathbb R^d)}
       +\sum_{j=1}^N\|E_{\eta_j}u\|_{W^{1,p}(\mathbb R^d)}\\
 &\le\left(C_{0,p}+\sum_{j=1}^N C_{\eta_j,F_j,p}\right)
       \|u\|_{W^{1,p}(\Omega)}.
 \end{split}
\]
常数包含有限个图表的 Lipschitz 界与截断函数的一阶导数界，这就证明\eqref{b1:eq:lipschitz-domain-extension-bound}。各项支集都包含在预先固定的紧集中，有限并仍紧。整个公式不随 \(p\) 改变，因此在不同指数空间的交上也给出同一个延拓。
\end{proof}

半空间的精确因子 \(2^{1/p}\) 来自标准坐标中的测度保持反射；一般区域还经过坐标变换和截断，不能期待相同常数。延拓算子也不是唯一的，改变外部截断或边界图表通常就会改变开集外的函数，但不会改变它在 \(\Omega\) 上的限制。

一阶构造不能直接重复用于高阶空间。取 \(\zeta\in C_c^\infty(\mathbb R)\)，使它在零点附近等于 \(1\)，并令 \(u(t)=t\zeta(t)\) 于 \(t>0\)。这个函数属于半轴上的每个整数阶 Sobolev 空间，但其偶反射在零点附近等于 \(|t|\)，二阶分布导数含有 \(2\delta_0\)，不是局部可积函数。另一方面，Lipschitz 图函数本身也没有足够的高阶导数供坐标链式法则使用。因此这里只证明一阶延拓；Lipschitz 区域上的高阶延拓需要另外的构造，不能把普通反射的失败解释为高阶延拓不可能。

对于有限 \(p\)，将 \(Eu\) 在整个空间磨光再限制回 \(\Omega\)，便可得到跨越边界的光滑近似；\(p=\infty\) 时，延拓的存在并不使卷积自动在一阶 Sobolev 范数下收敛。两个问题分别是能否补出受控函数、能否在指定范数中逼近，不能因前者成立而省略后者的端点条件。

% !TeX root = ../../Book1.tex
\section{局部化}
\label{b1:sec:2204}

前三节反复使用“先乘截断，再磨光或延拓”。现在把其中的估计单独整理出来。局部化不仅用来拼接函数，也能把一个只允许紧支集测试函数的积分恒等式，转成较小区域上的定量控制。需要留意的是截断函数的导数：它们记录了较小区域与外部边界之间的距离，也是局部估计中常数的来源。

\subsection{截断函数}
\label{b1:subsec:220401}

设 \(U,V\subseteq\Omega\) 为开集，且 \(\overline U\Subset V\)。由光滑截断引理，可选 \(0\le\eta\le1\)，使 \(\eta\in C_c^\infty(V)\)，并在 \(\overline U\) 的邻域内等于 \(1\)。对 \(u\in W^{k,p}_{\mathrm{loc}}(\Omega)\)，函数 \(\eta\cdot u\) 在 \(U\) 上与 \(u\) 相同，却具有内部紧支集。这个替换保留了所关注区域中的函数和全部弱导数。

\begin{proposition}\label{b1:prop:sobolev-cutoff-estimate}
设 \(k\in\mathbb N_0\)、\(1\le p\le\infty\)，且 \(\eta\in C_c^\infty(V)\)。对 \(u\in W^{k,p}(V)\)，把 \(\eta\cdot u\) 在 \(V\) 外补零，得到
\[
 \|\widetilde{\eta\cdot u}\|_{W^{k,p}(\mathbb R^d)}
 \le C_{d,k}
 \left(\max_{|\beta|\le k}\|\partial^\beta\eta\|_\infty\right)
 \|u\|_{W^{k,p}(V)}.
\]
更精细地，对每个 \(|\alpha|\le k\)，
\begin{equation}\label{b1:eq:sobolev-cutoff-components}
 \|\partial^\alpha(\eta\cdot u)\|_{L^p(V)}
 \le\sum_{\beta\le\alpha}\binom\alpha\beta
     \|\partial^\beta\eta\|_\infty
     \cdot\|\partial^{\alpha-\beta}u\|_{L^p(V)}.
\end{equation}
\end{proposition}
\begin{proof}
弱 Leibniz 公式给出
\[
 \partial^\alpha(\eta\cdot u)
 =\sum_{\beta\le\alpha}\binom\alpha\beta
       (\partial^\beta\eta)\cdot(\partial^{\alpha-\beta}u).
\]
对每一项使用有界函数乘法的 \(L^p\) 估计，再用三角不等式，即得分量估计。这对 \(p=\infty\) 也成立。多重指标的数目和二项式系数仅依赖于 \(d,k\)，有限维范数比较于是给出总范数估计；所用比较常数可统一取在 \(1\le p\le\infty\) 上有界。\cref{b1:lem:compact-sobolev-zero-extension}保证补零后的导数不增加边界项，也不改变范数。
\end{proof}

在两个同心球之间，截断常数可以写得更具体。若 \(0<r<R\)，存在 \(\eta\in C_c^\infty(B(x_0,R))\)，使 \(0\le\eta\le1\)、\(\eta=1\) 于 \(\overline B(x_0,r)\)，并且
\begin{equation}\label{b1:eq:ball-cutoff-derivatives}
 \|\partial^\beta\eta\|_\infty
 \le C_{d,\beta}(R-r)^{-|\beta|}
 \quad(|\beta|\ge1).
\end{equation}
例如取 \(\delta=(R-r)/4\)，将 \(B(x_0,r+2\delta)\) 的示性函数与 \(\rho_\delta\) 卷积。卷积在 \(\overline B(x_0,r)\) 附近等于 \(1\)，支集包含于 \(\overline B(x_0,r+3\delta)\subset B(x_0,R)\)。将导数放在核上，并用
\[
 \|\partial^\beta\rho_\delta\|_{L^1}
 =\delta^{-|\beta|}\|\partial^\beta\rho\|_{L^1},
\]
就得到上述界。

\ProseParagraphBreak

因此，局部化不是无代价的限制运算。例如一阶时，
\[
 \|\nabla(\eta\cdot u)\|_{L^p}
 \le \|\eta\cdot\nabla u\|_{L^p}
      +\|u\cdot\nabla\eta\|_{L^p}.
\]
第二项含有函数本身，而不只含梯度；当 \(R-r\) 缩小时，它通常会增大。把一个全域估计用于 \(\eta\cdot u\) 后，必须保留这项，不能因 \(\eta=1\) 于内球就将它删去。

\subsection{单位分解}
\label{b1:subsec:220402}

截断处理一个区域；单位分解则使相容的局部函数成为一个整体。这里先区分“局部属于 Sobolev 空间”和“具有有限全域范数”。

\begin{proposition}\label{b1:prop:sobolev-local-gluing}
设 \(k\in\mathbb N_0\)、\(1\le p\le\infty\)，且 \((U_j)_{j\in J}\) 为开集 \(\Omega\) 的开覆盖。若 \(u\in L^1_{\mathrm{loc}}(\Omega)\)，且每个限制 \(u|_{U_j}\) 都属于 \(W^{k,p}_{\mathrm{loc}}(U_j)\)，则 \(u\in W^{k,p}_{\mathrm{loc}}(\Omega)\)。

\ProseParagraphBreak

若进一步 \(u\in L^p(\Omega)\)，且各个局部弱导数拼成的函数 \(g_\alpha\) 满足 \(g_\alpha\in L^p(\Omega)\)（\(1\le|\alpha|\le k\)），则 \(u\in W^{k,p}(\Omega)\)，并有 \(\partial^\alpha u=g_\alpha\)。
\end{proposition}
\begin{proof}
在两个开集的交上，同阶弱导数由唯一性相等。因此各局部导数可以拼接。为避免不可数个零测集的问题，先利用欧氏空间的可数基，从原覆盖中取出可数子覆盖；再在这个可数族内去掉各交集上导数不一致的零测集。由此得到几乎处处定义且可测的 \(g_\alpha\)。

任取 \(K\Subset\Omega\)。可以用有限个相对紧于某些 \(U_j\) 的小开集覆盖 \(K\)，所以 \(g_\alpha\) 在 \(K\) 附近属于 \(L^p\)。现取 \(\varphi\in C_c^\infty(\Omega)\)，并将其按从属于覆盖的光滑单位分解写为有限和
\[
 \varphi=\sum_\ell\psi_\ell\varphi.
\]
每个非零分量都紧支于某个 \(U_j\)。在那里应用弱导数恒等式，求和得到
\[
 \begin{aligned}
 \int_\Omega u\,\partial^\alpha\varphi
 &=\sum_\ell\int_\Omega u\,\partial^\alpha(\psi_\ell\varphi)\\
 &=(-1)^{|\alpha|}\sum_\ell\int_\Omega g_\alpha\psi_\ell\varphi
 =(-1)^{|\alpha|}\int_\Omega g_\alpha\varphi.
 \end{aligned}
\]
第一处等号使用的是测试函数分解后的导数之和；单位分解导数产生的项已经在这个和中相消。于是 \(g_\alpha\) 是全域弱导数。局部或全域的可积性分别给出两种结论。
\end{proof}

同样，若存在 \(\Omega\) 中的局部有限闭集族 \((K_j)\)，使每个 \(v_j\in W^{k,p}_{\mathrm{loc}}(\Omega)\) 都在 \(\Omega\setminus K_j\) 上几乎处处为零，那么局部有限和
\[
 v=\sum_jv_j
\]
属于 \(W^{k,p}_{\mathrm{loc}}(\Omega)\)，其导数可以逐项相加，因为每个点附近实际只有有限项。这里没有保证 \(\|v\|_{W^{k,p}(\Omega)}<\infty\)。例如在全空间取光滑单位分解，所有分解函数都在各自的紧集外为零，但它们的和为 \(1\)，在有限 \(p\) 时不属于 \(L^p(\mathbb R^d)\)。

\ProseParagraphBreak

获得全域界需要额外估计。一种简单而充分的条件是
\[
 \sum_j\|v_j\|_{W^{k,p}(\Omega)}<\infty.
\]
此时完备性使级数在全域范数中收敛，其极限与局部有限和相同，而且总范数不超过右侧之和。Meyers--Serrin 定理正是对磨光误差建立了这样的可和界；不必要求单位分解后的原函数各项也满足它。

\subsection{局部 Sobolev 空间}
\label{b1:subsec:220403}

\secref{b1:sec:2101}的局部空间定义，要求在每个内部相对紧开集上控制函数及导数。用刚才的截断估计，可以把这个条件改写成全空间中的一族条件。

\begin{proposition}\label{b1:prop:local-sobolev-cutoff-characterization}
设 \(k\in\mathbb N_0\)、\(1\le p\le\infty\)，且 \(u\in L^1_{\mathrm{loc}}(\Omega)\)。下列条件等价：
\begin{equivalences}
\item\label{b1:item:local-sobolev-membership}
\(u\in W^{k,p}_{\mathrm{loc}}(\Omega)\)。
\item\label{b1:item:all-sobolev-cutoffs}
对每个 \(\chi\in C_c^\infty(\Omega)\)，补零函数 \(\widetilde{\chi\cdot u}\) 属于 \(W^{k,p}(\mathbb R^d)\)。
\end{equivalences}
若 \(u_n,u\in W^{k,p}_{\mathrm{loc}}(\Omega)\)，则在每个 \(\overline U\Subset\Omega\) 的开集 \(U\) 上有 \(u_n\to u\) 于 \(W^{k,p}(U)\)，当且仅当对每个上述 \(\chi\)，
\[
 \widetilde{\chi\cdot u_n}\longrightarrow\widetilde{\chi\cdot u}
 \quad\text{于 }W^{k,p}(\mathbb R^d).
\]
\end{proposition}
\begin{proof}
由局部可积性，先取一个内部开集包含 \(\operatorname{supp}\chi\)，再应用截断估计和补零引理，得到从\itemref{b1:item:local-sobolev-membership}到\itemref{b1:item:all-sobolev-cutoffs}的蕴含。反过来，对任意 \(\overline U\Subset\Omega\)，选取在 \(\overline U\) 附近为一的 \(\chi\)。由于 \(\widetilde{\chi\cdot u}|_U=u|_U\)，限制全空间弱导数便得到 \(u|_U\in W^{k,p}(U)\)。

收敛部分使用相同的两个操作。由局部收敛及固定 \(\chi\) 的截断估计，得到全空间乘积的收敛；反过来，取 \(\chi=1\) 于 \(\overline U\) 附近，限制不会增大范数，故得到 \(U\) 上的收敛。
\end{proof}

这里的“每个截断函数”不意味着实际工作中必须检验不可数多项。取递增的内部相对紧开集 \(U_j\) 穷竭 \(\Omega\)，再固定 \(\chi_j=1\) 于 \(\overline{U_j}\) 附近，即可只用这一列截断检验局部条件。任意内部紧集都包含在某个 \(U_j\) 中，因而不会漏掉局部信息。

\ProseParagraphBreak

局部有界性还足以在自反指数范围内抽取局部弱极限。具体说，若 \(1<p<\infty\)，且 \((u_n)\) 在每个 \(W^{k,p}(U_j)\) 中有界，则在 \(U_1\) 上用自反性抽取弱收敛子列，再在 \(U_2\) 上从中继续抽取，依次进行。对角子列在每个 \(U_j\) 上弱收敛。限制算子是有界线性的，故相邻区域上的极限限制相同；这些极限拼成 \(u\in W^{k,p}_{\mathrm{loc}}(\Omega)\)。这只给出局部弱收敛，不保证全域范数有界，更不保证强收敛。以互相远离的平移函数为例，函数可以在每个固定内部区域上消失，而全域范数一直保持不变。

\subsection{局部能量估计与 Caccioppoli 不等式}
\label{b1:subsec:220404}

局部化的一个重要用途，是允许测试函数依赖于待研究的 Sobolev 函数。若积分恒等式最初只对 \(C_c^\infty\) 函数成立，不能直接把一个不光滑的乘积代入；内部紧支集的范数逼近提供了所需过渡。

\ProseParagraphBreak

以下在实数范围内讨论。设 \(u\in H^1_{\mathrm{loc}}(\Omega)\)，满足
\begin{equation}\label{b1:eq:weak-harmonic-condition}
 \int_\Omega\nabla u\cdot\nabla\varphi\,\mathrm dx=0
 \quad\bigl(\varphi\in C_c^\infty(\Omega)\bigr).
\end{equation}
这个恒等式是方程 \(-\Delta u=0\) 的弱形式，并不预先要求 \(u\) 有经典二阶导数。下面的\term[caccioppoli-budengshi]{Caccioppoli 不等式}[Caccioppoli inequality]用较大球中的函数大小控制较小球中的梯度能量。

\begin{theorem}[Caccioppoli]\label{b1:thm:caccioppoli}
若 \(u\in H^1_{\mathrm{loc}}(\Omega)\) 满足\eqref{b1:eq:weak-harmonic-condition}，则对实常数 \(c\) 及实值 \(\eta\in C_c^\infty(\Omega)\)，
\begin{equation}\label{b1:eq:caccioppoli-cutoff}
 \int_\Omega\eta^2|\nabla u|^2\,\mathrm dx
 \le4\int_\Omega|u-c|^2\cdot|\nabla\eta|^2\,\mathrm dx.
\end{equation}
特别地，若 \(0<r<R\)、\(\overline B(x_0,R)\Subset\Omega\)，则
\begin{equation}\label{b1:eq:caccioppoli-balls}
 \int_{B(x_0,r)}|\nabla u|^2\,\mathrm dx
 \le\frac{C_d}{(R-r)^2}
       \int_{B(x_0,R)}|u-c|^2\,\mathrm dx.
\end{equation}
\end{theorem}
\begin{proof}
令 \(v=\eta^2(u-c)\)。弱乘积公式保证 \(v\in H^1(\Omega)\)，且它具有内部紧支集。选一个开集 \(V\)，使 \(\operatorname{supp}\eta\Subset V\)、\(\overline V\Subset\Omega\)。由\cref{b1:prop:compact-sobolev-smooth-approximation}，存在 \(v_j\in C_c^\infty(V)\)，使 \(v_j\to v\) 于 \(H^1(\Omega)\)。由于 \(\nabla u\in L^2(V)\)，Cauchy--Schwarz 不等式给出
\[
 \left|\int_V\nabla u\cdot\nabla(v_j-v)\,\mathrm dx\right|
 \le\|\nabla u\|_{L^2(V)}
      \cdot\|\nabla(v_j-v)\|_{L^2(V)}
 \longrightarrow0.
\]
因此弱恒等式也允许使用测试函数 \(v\)。展开
\[
 \nabla v=\eta^2\nabla u+2\eta(u-c)\nabla\eta
\]
并代入，得到
\[
 \int_\Omega\eta^2|\nabla u|^2
 =-2\int_\Omega\eta(u-c)\nabla u\cdot\nabla\eta.
\]
用向量 Cauchy--Schwarz 不等式及 \(2ab\le a^2/2+2b^2\)，得到
\[
 \int_\Omega\eta^2|\nabla u|^2
 \le\frac12\int_\Omega\eta^2|\nabla u|^2
       +2\int_\Omega|u-c|^2\cdot|\nabla\eta|^2.
\]
吸收左侧的一半即得\eqref{b1:eq:caccioppoli-cutoff}。最后选取\eqref{b1:eq:ball-cutoff-derivatives}中的球间截断，使用 \(\eta=1\) 于内球及 \(|\nabla\eta|\le C_d/(R-r)\)，便得\eqref{b1:eq:caccioppoli-balls}。
\end{proof}

常数 \(c\) 可以自由选取，因为方程只涉及梯度。取 \(c=0\) 得到直接的 \(L^2\) 控制；取 \(c\) 为较大球上的平均值，则右侧只测量函数相对于常数的振幅。当 \(u\) 本来就是常数时，左右两侧都为零，这也说明估计保留了方程对加常数的不变性。

\ProseParagraphBreak

本证明没有从 \(\|u\|_{L^2}\) 推出任意函数的梯度估计，而是先用方程将梯度能量变成截断交叉项。没有弱恒等式时，这一步并不存在。另一方面，截断只在内部起作用，所以定理没有使用边界值或迹；靠近边界的估计则需要额外信息。下一章将给出 Sobolev 函数的边界值应如何定义，以及这种边界值与 \(W_0^{1,p}\) 的关系。

单位分解的归一化与支撑安排见\cref{b1:fig:visual-f16}。
% F16：本书原创矢量示意；依据所在节的定义、证明或明确模型。
\begin{figure}[htbp]
\centering
\begin{tikzpicture}[x=1cm,y=1cm,>=stealth,every node/.style={font=\small},line width=.65pt]
\node at(6,3.7){开覆盖与截断函数};
\draw[AnalysisBlue,thick](1.0,2.7)--(5.0,2.7);\node[AnalysisBlue] at(3.0,3.1){\(U_1\)};
\draw[AnalysisBlue,thick] plot coordinates {(1.0000,0.0000) (1.0400,0.0000) (1.0800,0.0000) (1.1200,0.0174) (1.1600,0.0521) (1.2000,0.0868) (1.2400,0.1216) (1.2800,0.1563) (1.3200,0.1911) (1.3600,0.2258) (1.4000,0.2605) (1.4400,0.2953) (1.4800,0.3300) (1.5200,0.3647) (1.5600,0.3995) (1.6000,0.4342) (1.6400,0.4689) (1.6800,0.5037) (1.7200,0.5384) (1.7600,0.5732) (1.8000,0.6079) (1.8400,0.6426) (1.8800,0.6774) (1.9200,0.7121) (1.9600,0.7468) (2.0000,0.7816) (2.0400,0.8163) (2.0800,0.8511) (2.1200,0.8858) (2.1600,0.9205) (2.2000,0.9553) (2.2400,0.9900) (2.2800,1.0247) (2.3200,1.0595) (2.3600,1.0942) (2.4000,1.1289) (2.4400,1.1637) (2.4800,1.1984) (2.5200,1.2332) (2.5600,1.2679) (2.6000,1.3026) (2.6400,1.3374) (2.6800,1.3721) (2.7200,1.4068) (2.7600,1.4416) (2.8000,1.4763) (2.8400,1.5111) (2.8800,1.5458) (2.9200,1.5805) (2.9600,1.6153) (3.0000,1.6500) (3.0400,1.6153) (3.0800,1.5805) (3.1200,1.5458) (3.1600,1.5111) (3.2000,1.4763) (3.2400,1.4416) (3.2800,1.4068) (3.3200,1.3721) (3.3600,1.3374) (3.4000,1.3026) (3.4400,1.2679) (3.4800,1.2332) (3.5200,1.1984) (3.5600,1.1637) (3.6000,1.1289) (3.6400,1.0942) (3.6800,1.0595) (3.7200,1.0247) (3.7600,0.9900) (3.8000,0.9553) (3.8400,0.9205) (3.8800,0.8858) (3.9200,0.8511) (3.9600,0.8163) (4.0000,0.7816) (4.0400,0.7468) (4.0800,0.7121) (4.1200,0.6774) (4.1600,0.6426) (4.2000,0.6079) (4.2400,0.5732) (4.2800,0.5384) (4.3200,0.5037) (4.3600,0.4689) (4.4000,0.4342) (4.4400,0.3995) (4.4800,0.3647) (4.5200,0.3300) (4.5600,0.2953) (4.6000,0.2605) (4.6400,0.2258) (4.6800,0.1911) (4.7200,0.1563) (4.7600,0.1216) (4.8000,0.0868) (4.8400,0.0521) (4.8800,0.0174) (4.9200,0.0000) (4.9600,0.0000) (5.0000,0.0000)};
\draw[orange!80!black,thick](4.5,2.7)--(8.5,2.7);\node[orange!80!black] at(6.5,3.1){\(U_2\)};
\draw[orange!80!black,thick] plot coordinates {(4.5000,0.0000) (4.5400,0.0000) (4.5800,0.0000) (4.6200,0.0174) (4.6600,0.0521) (4.7000,0.0868) (4.7400,0.1216) (4.7800,0.1563) (4.8200,0.1911) (4.8600,0.2258) (4.9000,0.2605) (4.9400,0.2953) (4.9800,0.3300) (5.0200,0.3647) (5.0600,0.3995) (5.1000,0.4342) (5.1400,0.4689) (5.1800,0.5037) (5.2200,0.5384) (5.2600,0.5732) (5.3000,0.6079) (5.3400,0.6426) (5.3800,0.6774) (5.4200,0.7121) (5.4600,0.7468) (5.5000,0.7816) (5.5400,0.8163) (5.5800,0.8511) (5.6200,0.8858) (5.6600,0.9205) (5.7000,0.9553) (5.7400,0.9900) (5.7800,1.0247) (5.8200,1.0595) (5.8600,1.0942) (5.9000,1.1289) (5.9400,1.1637) (5.9800,1.1984) (6.0200,1.2332) (6.0600,1.2679) (6.1000,1.3026) (6.1400,1.3374) (6.1800,1.3721) (6.2200,1.4068) (6.2600,1.4416) (6.3000,1.4763) (6.3400,1.5111) (6.3800,1.5458) (6.4200,1.5805) (6.4600,1.6153) (6.5000,1.6500) (6.5400,1.6153) (6.5800,1.5805) (6.6200,1.5458) (6.6600,1.5111) (6.7000,1.4763) (6.7400,1.4416) (6.7800,1.4068) (6.8200,1.3721) (6.8600,1.3374) (6.9000,1.3026) (6.9400,1.2679) (6.9800,1.2332) (7.0200,1.1984) (7.0600,1.1637) (7.1000,1.1289) (7.1400,1.0942) (7.1800,1.0595) (7.2200,1.0247) (7.2600,0.9900) (7.3000,0.9553) (7.3400,0.9205) (7.3800,0.8858) (7.4200,0.8511) (7.4600,0.8163) (7.5000,0.7816) (7.5400,0.7468) (7.5800,0.7121) (7.6200,0.6774) (7.6600,0.6426) (7.7000,0.6079) (7.7400,0.5732) (7.7800,0.5384) (7.8200,0.5037) (7.8600,0.4689) (7.9000,0.4342) (7.9400,0.3995) (7.9800,0.3647) (8.0200,0.3300) (8.0600,0.2953) (8.1000,0.2605) (8.1400,0.2258) (8.1800,0.1911) (8.2200,0.1563) (8.2600,0.1216) (8.3000,0.0868) (8.3400,0.0521) (8.3800,0.0174) (8.4200,0.0000) (8.4600,0.0000) (8.5000,0.0000)};
\draw[black,thick](8.0,2.7)--(12.0,2.7);\node[black] at(10.0,3.1){\(U_3\)};
\draw[black,thick] plot coordinates {(8.0000,0.0000) (8.0400,0.0000) (8.0800,0.0000) (8.1200,0.0174) (8.1600,0.0521) (8.2000,0.0868) (8.2400,0.1216) (8.2800,0.1563) (8.3200,0.1911) (8.3600,0.2258) (8.4000,0.2605) (8.4400,0.2953) (8.4800,0.3300) (8.5200,0.3647) (8.5600,0.3995) (8.6000,0.4342) (8.6400,0.4689) (8.6800,0.5037) (8.7200,0.5384) (8.7600,0.5732) (8.8000,0.6079) (8.8400,0.6426) (8.8800,0.6774) (8.9200,0.7121) (8.9600,0.7468) (9.0000,0.7816) (9.0400,0.8163) (9.0800,0.8511) (9.1200,0.8858) (9.1600,0.9205) (9.2000,0.9553) (9.2400,0.9900) (9.2800,1.0247) (9.3200,1.0595) (9.3600,1.0942) (9.4000,1.1289) (9.4400,1.1637) (9.4800,1.1984) (9.5200,1.2332) (9.5600,1.2679) (9.6000,1.3026) (9.6400,1.3374) (9.6800,1.3721) (9.7200,1.4068) (9.7600,1.4416) (9.8000,1.4763) (9.8400,1.5111) (9.8800,1.5458) (9.9200,1.5805) (9.9600,1.6153) (10.0000,1.6500) (10.0400,1.6153) (10.0800,1.5805) (10.1200,1.5458) (10.1600,1.5111) (10.2000,1.4763) (10.2400,1.4416) (10.2800,1.4068) (10.3200,1.3721) (10.3600,1.3374) (10.4000,1.3026) (10.4400,1.2679) (10.4800,1.2332) (10.5200,1.1984) (10.5600,1.1637) (10.6000,1.1289) (10.6400,1.0942) (10.6800,1.0595) (10.7200,1.0247) (10.7600,0.9900) (10.8000,0.9553) (10.8400,0.9205) (10.8800,0.8858) (10.9200,0.8511) (10.9600,0.8163) (11.0000,0.7816) (11.0400,0.7468) (11.0800,0.7121) (11.1200,0.6774) (11.1600,0.6426) (11.2000,0.6079) (11.2400,0.5732) (11.2800,0.5384) (11.3200,0.5037) (11.3600,0.4689) (11.4000,0.4342) (11.4400,0.3995) (11.4800,0.3647) (11.5200,0.3300) (11.5600,0.2953) (11.6000,0.2605) (11.6400,0.2258) (11.6800,0.1911) (11.7200,0.1563) (11.7600,0.1216) (11.8000,0.0868) (11.8400,0.0521) (11.8800,0.0174) (11.9200,0.0000) (11.9600,0.0000) (12.0000,0.0000)};
\node at(6,-.65){\(\psi_j=\eta_j/\sum_i\eta_i\)，在工作区域上 \(\sum_j\psi_j=1\)};
\draw[AnalysisBlue,thick] plot coordinates {(1.2000,-1.3000) (1.2463,-1.3000) (1.2926,-1.3000) (1.3389,-1.3000) (1.3852,-1.3000) (1.4314,-1.3000) (1.4777,-1.3000) (1.5240,-1.3000) (1.5703,-1.3000) (1.6166,-1.3000) (1.6629,-1.3000) (1.7092,-1.3000) (1.7555,-1.3000) (1.8017,-1.3000) (1.8480,-1.3000) (1.8943,-1.3000) (1.9406,-1.3000) (1.9869,-1.3000) (2.0332,-1.3000) (2.0795,-1.3000) (2.1258,-1.3000) (2.1721,-1.3000) (2.2183,-1.3000) (2.2646,-1.3000) (2.3109,-1.3000) (2.3572,-1.3000) (2.4035,-1.3000) (2.4498,-1.3000) (2.4961,-1.3000) (2.5424,-1.3000) (2.5886,-1.3000) (2.6349,-1.3000) (2.6812,-1.3000) (2.7275,-1.3000) (2.7738,-1.3000) (2.8201,-1.3000) (2.8664,-1.3000) (2.9127,-1.3000) (2.9590,-1.3000) (3.0052,-1.3000) (3.0515,-1.3000) (3.0978,-1.3000) (3.1441,-1.3000) (3.1904,-1.3000) (3.2367,-1.3000) (3.2830,-1.3000) (3.3293,-1.3000) (3.3755,-1.3000) (3.4218,-1.3000) (3.4681,-1.3000) (3.5144,-1.3000) (3.5607,-1.3000) (3.6070,-1.3000) (3.6533,-1.3000) (3.6996,-1.3000) (3.7459,-1.3000) (3.7921,-1.3000) (3.8384,-1.3000) (3.8847,-1.3000) (3.9310,-1.3000) (3.9773,-1.3000) (4.0236,-1.3000) (4.0699,-1.3000) (4.1162,-1.3000) (4.1624,-1.3000) (4.2087,-1.3000) (4.2550,-1.3000) (4.3013,-1.3000) (4.3476,-1.3000) (4.3939,-1.3000) (4.4402,-1.3000) (4.4865,-1.3000) (4.5328,-1.3000) (4.5790,-1.3000) (4.6253,-1.4266) (4.6716,-1.6581) (4.7179,-1.8895) (4.7642,-2.1210) (4.8105,-2.3524) (4.8568,-2.5838) (4.9031,-2.8000) (4.9493,-2.8000) (4.9956,-2.8000) (5.0419,-2.8000) (5.0882,-2.8000) (5.1345,-2.8000) (5.1808,-2.8000) (5.2271,-2.8000) (5.2734,-2.8000) (5.3197,-2.8000) (5.3659,-2.8000) (5.4122,-2.8000) (5.4585,-2.8000) (5.5048,-2.8000) (5.5511,-2.8000) (5.5974,-2.8000) (5.6437,-2.8000) (5.6900,-2.8000) (5.7362,-2.8000) (5.7825,-2.8000) (5.8288,-2.8000) (5.8751,-2.8000) (5.9214,-2.8000) (5.9677,-2.8000) (6.0140,-2.8000) (6.0603,-2.8000) (6.1066,-2.8000) (6.1528,-2.8000) (6.1991,-2.8000) (6.2454,-2.8000) (6.2917,-2.8000) (6.3380,-2.8000) (6.3843,-2.8000) (6.4306,-2.8000) (6.4769,-2.8000) (6.5231,-2.8000) (6.5694,-2.8000) (6.6157,-2.8000) (6.6620,-2.8000) (6.7083,-2.8000) (6.7546,-2.8000) (6.8009,-2.8000) (6.8472,-2.8000) (6.8934,-2.8000) (6.9397,-2.8000) (6.9860,-2.8000) (7.0323,-2.8000) (7.0786,-2.8000) (7.1249,-2.8000) (7.1712,-2.8000) (7.2175,-2.8000) (7.2638,-2.8000) (7.3100,-2.8000) (7.3563,-2.8000) (7.4026,-2.8000) (7.4489,-2.8000) (7.4952,-2.8000) (7.5415,-2.8000) (7.5878,-2.8000) (7.6341,-2.8000) (7.6803,-2.8000) (7.7266,-2.8000) (7.7729,-2.8000) (7.8192,-2.8000) (7.8655,-2.8000) (7.9118,-2.8000) (7.9581,-2.8000) (8.0044,-2.8000) (8.0507,-2.8000) (8.0969,-2.8000) (8.1432,-2.8000) (8.1895,-2.8000) (8.2358,-2.8000) (8.2821,-2.8000) (8.3284,-2.8000) (8.3747,-2.8000) (8.4210,-2.8000) (8.4672,-2.8000) (8.5135,-2.8000) (8.5598,-2.8000) (8.6061,-2.8000) (8.6524,-2.8000) (8.6987,-2.8000) (8.7450,-2.8000) (8.7913,-2.8000) (8.8376,-2.8000) (8.8838,-2.8000) (8.9301,-2.8000) (8.9764,-2.8000) (9.0227,-2.8000) (9.0690,-2.8000) (9.1153,-2.8000) (9.1616,-2.8000) (9.2079,-2.8000) (9.2541,-2.8000) (9.3004,-2.8000) (9.3467,-2.8000) (9.3930,-2.8000) (9.4393,-2.8000) (9.4856,-2.8000) (9.5319,-2.8000) (9.5782,-2.8000) (9.6245,-2.8000) (9.6707,-2.8000) (9.7170,-2.8000) (9.7633,-2.8000) (9.8096,-2.8000) (9.8559,-2.8000) (9.9022,-2.8000) (9.9485,-2.8000) (9.9948,-2.8000) (10.0410,-2.8000) (10.0873,-2.8000) (10.1336,-2.8000) (10.1799,-2.8000) (10.2262,-2.8000) (10.2725,-2.8000) (10.3188,-2.8000) (10.3651,-2.8000) (10.4114,-2.8000) (10.4576,-2.8000) (10.5039,-2.8000) (10.5502,-2.8000) (10.5965,-2.8000) (10.6428,-2.8000) (10.6891,-2.8000) (10.7354,-2.8000) (10.7817,-2.8000) (10.8279,-2.8000) (10.8742,-2.8000) (10.9205,-2.8000) (10.9668,-2.8000) (11.0131,-2.8000) (11.0594,-2.8000) (11.1057,-2.8000) (11.1520,-2.8000) (11.1983,-2.8000) (11.2445,-2.8000) (11.2908,-2.8000) (11.3371,-2.8000) (11.3834,-2.8000) (11.4297,-2.8000) (11.4760,-2.8000) (11.5223,-2.8000) (11.5686,-2.8000) (11.6148,-2.8000) (11.6611,-2.8000) (11.7074,-2.8000) (11.7537,-2.8000) (11.8000,-2.8000)};
\draw[orange!80!black,thick] plot coordinates {(1.2000,-2.8000) (1.2463,-2.8000) (1.2926,-2.8000) (1.3389,-2.8000) (1.3852,-2.8000) (1.4314,-2.8000) (1.4777,-2.8000) (1.5240,-2.8000) (1.5703,-2.8000) (1.6166,-2.8000) (1.6629,-2.8000) (1.7092,-2.8000) (1.7555,-2.8000) (1.8017,-2.8000) (1.8480,-2.8000) (1.8943,-2.8000) (1.9406,-2.8000) (1.9869,-2.8000) (2.0332,-2.8000) (2.0795,-2.8000) (2.1258,-2.8000) (2.1721,-2.8000) (2.2183,-2.8000) (2.2646,-2.8000) (2.3109,-2.8000) (2.3572,-2.8000) (2.4035,-2.8000) (2.4498,-2.8000) (2.4961,-2.8000) (2.5424,-2.8000) (2.5886,-2.8000) (2.6349,-2.8000) (2.6812,-2.8000) (2.7275,-2.8000) (2.7738,-2.8000) (2.8201,-2.8000) (2.8664,-2.8000) (2.9127,-2.8000) (2.9590,-2.8000) (3.0052,-2.8000) (3.0515,-2.8000) (3.0978,-2.8000) (3.1441,-2.8000) (3.1904,-2.8000) (3.2367,-2.8000) (3.2830,-2.8000) (3.3293,-2.8000) (3.3755,-2.8000) (3.4218,-2.8000) (3.4681,-2.8000) (3.5144,-2.8000) (3.5607,-2.8000) (3.6070,-2.8000) (3.6533,-2.8000) (3.6996,-2.8000) (3.7459,-2.8000) (3.7921,-2.8000) (3.8384,-2.8000) (3.8847,-2.8000) (3.9310,-2.8000) (3.9773,-2.8000) (4.0236,-2.8000) (4.0699,-2.8000) (4.1162,-2.8000) (4.1624,-2.8000) (4.2087,-2.8000) (4.2550,-2.8000) (4.3013,-2.8000) (4.3476,-2.8000) (4.3939,-2.8000) (4.4402,-2.8000) (4.4865,-2.8000) (4.5328,-2.8000) (4.5790,-2.8000) (4.6253,-2.6734) (4.6716,-2.4419) (4.7179,-2.2105) (4.7642,-1.9790) (4.8105,-1.7476) (4.8568,-1.5162) (4.9031,-1.3000) (4.9493,-1.3000) (4.9956,-1.3000) (5.0419,-1.3000) (5.0882,-1.3000) (5.1345,-1.3000) (5.1808,-1.3000) (5.2271,-1.3000) (5.2734,-1.3000) (5.3197,-1.3000) (5.3659,-1.3000) (5.4122,-1.3000) (5.4585,-1.3000) (5.5048,-1.3000) (5.5511,-1.3000) (5.5974,-1.3000) (5.6437,-1.3000) (5.6900,-1.3000) (5.7362,-1.3000) (5.7825,-1.3000) (5.8288,-1.3000) (5.8751,-1.3000) (5.9214,-1.3000) (5.9677,-1.3000) (6.0140,-1.3000) (6.0603,-1.3000) (6.1066,-1.3000) (6.1528,-1.3000) (6.1991,-1.3000) (6.2454,-1.3000) (6.2917,-1.3000) (6.3380,-1.3000) (6.3843,-1.3000) (6.4306,-1.3000) (6.4769,-1.3000) (6.5231,-1.3000) (6.5694,-1.3000) (6.6157,-1.3000) (6.6620,-1.3000) (6.7083,-1.3000) (6.7546,-1.3000) (6.8009,-1.3000) (6.8472,-1.3000) (6.8934,-1.3000) (6.9397,-1.3000) (6.9860,-1.3000) (7.0323,-1.3000) (7.0786,-1.3000) (7.1249,-1.3000) (7.1712,-1.3000) (7.2175,-1.3000) (7.2638,-1.3000) (7.3100,-1.3000) (7.3563,-1.3000) (7.4026,-1.3000) (7.4489,-1.3000) (7.4952,-1.3000) (7.5415,-1.3000) (7.5878,-1.3000) (7.6341,-1.3000) (7.6803,-1.3000) (7.7266,-1.3000) (7.7729,-1.3000) (7.8192,-1.3000) (7.8655,-1.3000) (7.9118,-1.3000) (7.9581,-1.3000) (8.0044,-1.3000) (8.0507,-1.3000) (8.0969,-1.3000) (8.1432,-1.5162) (8.1895,-1.7476) (8.2358,-1.9790) (8.2821,-2.2105) (8.3284,-2.4419) (8.3747,-2.6734) (8.4210,-2.8000) (8.4672,-2.8000) (8.5135,-2.8000) (8.5598,-2.8000) (8.6061,-2.8000) (8.6524,-2.8000) (8.6987,-2.8000) (8.7450,-2.8000) (8.7913,-2.8000) (8.8376,-2.8000) (8.8838,-2.8000) (8.9301,-2.8000) (8.9764,-2.8000) (9.0227,-2.8000) (9.0690,-2.8000) (9.1153,-2.8000) (9.1616,-2.8000) (9.2079,-2.8000) (9.2541,-2.8000) (9.3004,-2.8000) (9.3467,-2.8000) (9.3930,-2.8000) (9.4393,-2.8000) (9.4856,-2.8000) (9.5319,-2.8000) (9.5782,-2.8000) (9.6245,-2.8000) (9.6707,-2.8000) (9.7170,-2.8000) (9.7633,-2.8000) (9.8096,-2.8000) (9.8559,-2.8000) (9.9022,-2.8000) (9.9485,-2.8000) (9.9948,-2.8000) (10.0410,-2.8000) (10.0873,-2.8000) (10.1336,-2.8000) (10.1799,-2.8000) (10.2262,-2.8000) (10.2725,-2.8000) (10.3188,-2.8000) (10.3651,-2.8000) (10.4114,-2.8000) (10.4576,-2.8000) (10.5039,-2.8000) (10.5502,-2.8000) (10.5965,-2.8000) (10.6428,-2.8000) (10.6891,-2.8000) (10.7354,-2.8000) (10.7817,-2.8000) (10.8279,-2.8000) (10.8742,-2.8000) (10.9205,-2.8000) (10.9668,-2.8000) (11.0131,-2.8000) (11.0594,-2.8000) (11.1057,-2.8000) (11.1520,-2.8000) (11.1983,-2.8000) (11.2445,-2.8000) (11.2908,-2.8000) (11.3371,-2.8000) (11.3834,-2.8000) (11.4297,-2.8000) (11.4760,-2.8000) (11.5223,-2.8000) (11.5686,-2.8000) (11.6148,-2.8000) (11.6611,-2.8000) (11.7074,-2.8000) (11.7537,-2.8000) (11.8000,-2.8000)};
\draw[black,thick] plot coordinates {(1.2000,-2.8000) (1.2463,-2.8000) (1.2926,-2.8000) (1.3389,-2.8000) (1.3852,-2.8000) (1.4314,-2.8000) (1.4777,-2.8000) (1.5240,-2.8000) (1.5703,-2.8000) (1.6166,-2.8000) (1.6629,-2.8000) (1.7092,-2.8000) (1.7555,-2.8000) (1.8017,-2.8000) (1.8480,-2.8000) (1.8943,-2.8000) (1.9406,-2.8000) (1.9869,-2.8000) (2.0332,-2.8000) (2.0795,-2.8000) (2.1258,-2.8000) (2.1721,-2.8000) (2.2183,-2.8000) (2.2646,-2.8000) (2.3109,-2.8000) (2.3572,-2.8000) (2.4035,-2.8000) (2.4498,-2.8000) (2.4961,-2.8000) (2.5424,-2.8000) (2.5886,-2.8000) (2.6349,-2.8000) (2.6812,-2.8000) (2.7275,-2.8000) (2.7738,-2.8000) (2.8201,-2.8000) (2.8664,-2.8000) (2.9127,-2.8000) (2.9590,-2.8000) (3.0052,-2.8000) (3.0515,-2.8000) (3.0978,-2.8000) (3.1441,-2.8000) (3.1904,-2.8000) (3.2367,-2.8000) (3.2830,-2.8000) (3.3293,-2.8000) (3.3755,-2.8000) (3.4218,-2.8000) (3.4681,-2.8000) (3.5144,-2.8000) (3.5607,-2.8000) (3.6070,-2.8000) (3.6533,-2.8000) (3.6996,-2.8000) (3.7459,-2.8000) (3.7921,-2.8000) (3.8384,-2.8000) (3.8847,-2.8000) (3.9310,-2.8000) (3.9773,-2.8000) (4.0236,-2.8000) (4.0699,-2.8000) (4.1162,-2.8000) (4.1624,-2.8000) (4.2087,-2.8000) (4.2550,-2.8000) (4.3013,-2.8000) (4.3476,-2.8000) (4.3939,-2.8000) (4.4402,-2.8000) (4.4865,-2.8000) (4.5328,-2.8000) (4.5790,-2.8000) (4.6253,-2.8000) (4.6716,-2.8000) (4.7179,-2.8000) (4.7642,-2.8000) (4.8105,-2.8000) (4.8568,-2.8000) (4.9031,-2.8000) (4.9493,-2.8000) (4.9956,-2.8000) (5.0419,-2.8000) (5.0882,-2.8000) (5.1345,-2.8000) (5.1808,-2.8000) (5.2271,-2.8000) (5.2734,-2.8000) (5.3197,-2.8000) (5.3659,-2.8000) (5.4122,-2.8000) (5.4585,-2.8000) (5.5048,-2.8000) (5.5511,-2.8000) (5.5974,-2.8000) (5.6437,-2.8000) (5.6900,-2.8000) (5.7362,-2.8000) (5.7825,-2.8000) (5.8288,-2.8000) (5.8751,-2.8000) (5.9214,-2.8000) (5.9677,-2.8000) (6.0140,-2.8000) (6.0603,-2.8000) (6.1066,-2.8000) (6.1528,-2.8000) (6.1991,-2.8000) (6.2454,-2.8000) (6.2917,-2.8000) (6.3380,-2.8000) (6.3843,-2.8000) (6.4306,-2.8000) (6.4769,-2.8000) (6.5231,-2.8000) (6.5694,-2.8000) (6.6157,-2.8000) (6.6620,-2.8000) (6.7083,-2.8000) (6.7546,-2.8000) (6.8009,-2.8000) (6.8472,-2.8000) (6.8934,-2.8000) (6.9397,-2.8000) (6.9860,-2.8000) (7.0323,-2.8000) (7.0786,-2.8000) (7.1249,-2.8000) (7.1712,-2.8000) (7.2175,-2.8000) (7.2638,-2.8000) (7.3100,-2.8000) (7.3563,-2.8000) (7.4026,-2.8000) (7.4489,-2.8000) (7.4952,-2.8000) (7.5415,-2.8000) (7.5878,-2.8000) (7.6341,-2.8000) (7.6803,-2.8000) (7.7266,-2.8000) (7.7729,-2.8000) (7.8192,-2.8000) (7.8655,-2.8000) (7.9118,-2.8000) (7.9581,-2.8000) (8.0044,-2.8000) (8.0507,-2.8000) (8.0969,-2.8000) (8.1432,-2.5838) (8.1895,-2.3524) (8.2358,-2.1210) (8.2821,-1.8895) (8.3284,-1.6581) (8.3747,-1.4266) (8.4210,-1.3000) (8.4672,-1.3000) (8.5135,-1.3000) (8.5598,-1.3000) (8.6061,-1.3000) (8.6524,-1.3000) (8.6987,-1.3000) (8.7450,-1.3000) (8.7913,-1.3000) (8.8376,-1.3000) (8.8838,-1.3000) (8.9301,-1.3000) (8.9764,-1.3000) (9.0227,-1.3000) (9.0690,-1.3000) (9.1153,-1.3000) (9.1616,-1.3000) (9.2079,-1.3000) (9.2541,-1.3000) (9.3004,-1.3000) (9.3467,-1.3000) (9.3930,-1.3000) (9.4393,-1.3000) (9.4856,-1.3000) (9.5319,-1.3000) (9.5782,-1.3000) (9.6245,-1.3000) (9.6707,-1.3000) (9.7170,-1.3000) (9.7633,-1.3000) (9.8096,-1.3000) (9.8559,-1.3000) (9.9022,-1.3000) (9.9485,-1.3000) (9.9948,-1.3000) (10.0410,-1.3000) (10.0873,-1.3000) (10.1336,-1.3000) (10.1799,-1.3000) (10.2262,-1.3000) (10.2725,-1.3000) (10.3188,-1.3000) (10.3651,-1.3000) (10.4114,-1.3000) (10.4576,-1.3000) (10.5039,-1.3000) (10.5502,-1.3000) (10.5965,-1.3000) (10.6428,-1.3000) (10.6891,-1.3000) (10.7354,-1.3000) (10.7817,-1.3000) (10.8279,-1.3000) (10.8742,-1.3000) (10.9205,-1.3000) (10.9668,-1.3000) (11.0131,-1.3000) (11.0594,-1.3000) (11.1057,-1.3000) (11.1520,-1.3000) (11.1983,-1.3000) (11.2445,-1.3000) (11.2908,-1.3000) (11.3371,-1.3000) (11.3834,-1.3000) (11.4297,-1.3000) (11.4760,-1.3000) (11.5223,-1.3000) (11.5686,-1.3000) (11.6148,-1.3000) (11.6611,-1.3000) (11.7074,-1.3000) (11.7537,-1.3000) (11.8000,-1.3000)};
\draw[gray,dashed](1.2,-1.3)--(11.8,-1.3);\node[anchor=east] at(1,-1.3){\(1\)};\node[anchor=east] at(1,-2.8){\(0\)};\node at(6,-3.3){归一化后的三个权重};\end{tikzpicture}
\caption[局部覆盖与单位分解拼接]{局部覆盖与单位分解拼接。折线仅示意非负截断的支撑与重叠，实际选取光滑截断。分母在工作区域上为正，归一化后权重和为一；可数情形还需局部有限性。}
\label{b1:fig:visual-f16}
\end{figure}



\begin{chaptersummary}
内部紧支集使补零与弱微分相容；有限指数的卷积逼近随后给出光滑替代。任意开集上的 Meyers--Serrin 定理还需要逐层控制支集及误差，不能把单位分解的局部有限性误当作全域范数收敛。

而 \(W_0^{1,p}\) 用范数闭包定义，不依赖预先存在的边界值。它的元素可以补零，但任意开集上不能由补零反推出闭包条件。有限区间的零端点刻画说明了边界层估计如何起作用，也显示无穷指数与有限指数不能混用。

Lipschitz 区域的一阶延拓来自相容的反射、坐标变换和有限拼接。局部化则把这些构造中反复出现的截断整理为估计工具，使弱恒等式可以接受内部紧支集的 Sobolev 测试函数。下一章将正式定义迹，并说明在适当边界上，范数闭包意义的零边界条件怎样成为零迹条件。
\end{chaptersummary}

\BookExercises

% \chapref{b1:ch:22}习题临时稿；由章聚合文件已有的 BookExercises 入口统一接入。
% 十题均为自拟的基础训练或工具变式，来源说明记录实际前置，不冒认教材原题编号。
% 不使用第23章的迹或容量，也不使用后续嵌入及一般椭圆方程理论。

\begin{exercise}\label{b1:ex:22-basic-sobolev-mollification}
设 \(1\le p<\infty\)、\(u\in W^{1,p}(\mathbb R^d)\)，并令
\(u_\varepsilon=\rho_\varepsilon*u\)。逐个导数分量证明
\[
 \partial_i u_\varepsilon=\rho_\varepsilon*\partial_i u,qquad
 u_\varepsilon\longrightarrow u\quad\text{于 }W^{1,p}(\mathbb R^d).
\]
说明为什么同一论证不能把收敛结论直接扩展到 \(p=\infty\)。
\end{exercise}
% 来源：自拟；把正文全阶磨光命题降到一阶逐分量复算，训练最基本的Sobolev范数逼近。
\begin{solution}
由分布卷积的求导公式，
\[
 \partial_i u_\varepsilon
 =\partial_i(\rho_\varepsilon*u)
 =\rho_\varepsilon*\partial_i u.
\]
有限指数的逼近恒等分别给出
\[
 \|\rho_\varepsilon*u-u\|_p\longrightarrow0,qquad
 \|\rho_\varepsilon*\partial_i u-\partial_i u\|_p\longrightarrow0
 \quad(1\le i\le d).
\]
按本书的 \(W^{1,p}\) 范数把这有限多个分量的 \(p\) 次方相加，便得
\(\|u_\varepsilon-u\|_{W^{1,p}}\to0\)。无穷指数下仍有卷积范数不增，但一般 \(L^\infty\) 函数不具平移连续性；即使函数本身可以一致逼近，一阶导数的跳跃也可能使 \(W^{1,\infty}\) 误差不趋零。因此上述最后一步所用的有限指数逼近恒等不能照搬。
\end{solution}

\begin{exercise}\label{b1:ex:22-basic-even-reflection}
在半直线 \((0,\infty)\) 上取 \(u(t)=e^{-t}\)，并在全实轴定义偶反射
\(Eu(t)=e^{-|t|}\)。直接按测试函数定义证明
\[
 (Eu)'(t)=-\operatorname{sgn}(t)e^{-|t|}
 \quad\text{几乎处处},
\]
特别说明原点处没有 Dirac 项。再验证有限 \(p\) 时
\(\|Eu\|_{W^{1,p}(\mathbb R)}=2^{1/p}\|u\|_{W^{1,p}(0,\infty)}\)，而 \(p=\infty\) 时两侧相等。
\end{exercise}
% 来源：自拟；用显式偶反射训练跨边界弱求导及范数计算，不重复一般维半空间证明。
\begin{solution}
任取 \(\varphi\in C_c^\infty(\mathbb R)\)。分别在两侧积分并换元，得到
\[
 \begin{aligned}
 \int_{\mathbb R}e^{-|t|}\varphi'(t)\,\mathrm dt
 &=\int_0^\infty e^{-t}\varphi'(t)\,\mathrm dt
   +\int_0^\infty e^{-t}\varphi'(-t)\,\mathrm dt\\
 &=-\int_0^\infty(-e^{-t})\varphi(t)\,\mathrm dt
   -\int_{-\infty}^0 e^t\varphi(t)\,\mathrm dt.
 \end{aligned}
\]
两次分部积分在原点产生的边界项互相抵消；这正是偶反射两侧函数值相接所排除的 Dirac 项。右端等于
\(-\int_{\mathbb R}[-\operatorname{sgn}(t)e^{-|t|}]\varphi(t)\,\mathrm dt\)，故得到所列弱导数。

函数及弱导数的绝对值都等于 \(e^{-|t|}\)。有限 \(p\) 时，它们各自在全轴上的 \(p\) 次积分都是半轴上的两倍，所以 Sobolev 范数乘以 \(2^{1/p}\)；无穷端点的本性上确界在反射后不变。
\end{solution}

\begin{exercise}\label{b1:ex:22-strip-poincare}
设非空开集 \(\Omega\subset(a,b)\times\mathbb R^{d-1}\)，其中 \(L=b-a<\infty\)。证明对每个 \(1\le p\le\infty\) 及 \(u\in W_0^{1,p}(\Omega)\)，
\[
 \|u\|_{L^p(\Omega)}\le L\|\partial_1u\|_{L^p(\Omega)}.
\]
据此说明，仅由全部一阶弱导数组成的半范数在 \(W_0^{1,p}(\Omega)\) 上是范数，并与原 Sobolev 范数等价。不要求 \(\Omega\) 在其余坐标方向有界。
\end{exercise}
% 来源：自拟；紧支集光滑函数的沿线积分、Hölder及闭包定义给出粗Poincare估计，不预借后续同名定理。
\begin{solution}
先取 \(\varphi\in C_c^\infty(\Omega)\)，并把它补零为全空间上的紧支集光滑函数。写 \(x=(x_1,x')\)。沿第一坐标应用微积分基本定理，
\[
 \varphi(x_1,x')=\int_a^{x_1}\partial_1\varphi(t,x')\,\mathrm dt
 \quad(a<x_1<b).
\]
若 \(1\le p<\infty\)，Hölder 不等式给出
\[
 |\varphi(x_1,x')|^p
 \le L^{p-1}\int_a^b|\partial_1\varphi(t,x')|^p\,\mathrm dt.
\]
这里 \(p=1\) 时直接使用积分的三角不等式。先对 \(x_1\in(a,b)\) 积分，再对 \(x'\) 积分；补零函数及其导数在 \(\Omega\) 外都为零，所以
\[
 \|\varphi\|_{L^p(\Omega)}^p
 \le L^p\|\partial_1\varphi\|_{L^p(\Omega)}^p.
\]
若 \(p=\infty\)，直接由沿线积分取得
\(\|\varphi\|_\infty\le L\|\partial_1\varphi\|_\infty\)。

对一般 \(u\in W_0^{1,p}(\Omega)\)，按\cref{b1:def:sobolev-zero-boundary}取 \(\varphi_n\in C_c^\infty(\Omega)\)，使 \(\varphi_n\to u\) 于 \(W^{1,p}\)。函数及第一偏导数的范数均随之收敛，在刚证的不等式中取极限即可。这个步骤也适用于 \(p=\infty\)，因为使用的是闭包定义给定的逼近，而非一般函数的无穷范数磨光。

记 \(M_p(u)=\bigl(\sum_{j=1}^d\|\partial_ju\|_p^p\bigr)^{1/p}\)，其中 \(p<\infty\)；无穷端点改记 \(M_\infty(u)=\max_j\|\partial_ju\|_\infty\)。如果 \(M_p(u)=0\)，所求不等式使 \(\|u\|_p=0\)，故它确为范数。有限 \(p\) 时，
\[
 M_p(u)\le\|u\|_{W^{1,p}}
 \le(1+L^p)^{1/p}M_p(u);
\]
无穷端点则有
\[
 M_\infty(u)\le\|u\|_{W^{1,\infty}}
 \le\max(1,L)M_\infty(u).
\]
这些常数只取决于条带宽度和指数，并不要求条带中每条截线与 \(\Omega\) 的交集连通。
\end{solution}

\begin{exercise}[\(\star\) 选读]\label{b1:ex:22-w0-infty-interval}
设 \(I=(0,1)\)，\(u\in W^{1,\infty}(I)\)。证明下列两条等价：
\begin{equivalences}
\item\label{b1:item:22-w0-infty-closure}
\(u\in W_0^{1,\infty}(I)\)。
\item\label{b1:item:22-w0-infty-c1-endpoints}
等价类 \(u\) 有一个代表 \(U\in C^1([0,1])\)，且
\[
 U(0)=U(1)=U'(0)=U'(1)=0.
\]
\end{equivalences}
这里端点导数按单侧导数解释。
\end{exercise}
% 来源：自拟；用C1一致极限刻画无穷端点的闭包，进一步说明正文单个端点反例背后的完整条件。
\begin{solution}
先设 \itemref{b1:item:22-w0-infty-closure}成立，选取紧支集光滑函数 \(\varphi_n\to u\) 于 \(W^{1,\infty}(I)\)。每个 \(\varphi_n\) 在端点附近为零，可以看作 \([0,1]\) 上的光滑函数。连续函数在这个闭区间上的上确界等于其在 \(I\) 上的本性上确界。因此 \((\varphi_n)\) 与 \((\varphi_n')\) 都在一致范数下为 Cauchy 列，分别趋于连续函数 \(U,G\)，并满足
\[
 U(0)=U(1)=G(0)=G(1)=0.
\]
在等式 \(\varphi_n(x)=\int_0^x\varphi_n'(t)\,\mathrm dt\) 中取一致极限，得到 \(U(x)=\int_0^xG(t)\,\mathrm dt\)。所以 \(U\in C^1([0,1])\)，且 \(U'=G\)。由于 \(\varphi_n\) 同时在 \(L^\infty(I)\) 中趋于 \(u\)，这个 \(U\) 正是 \(u\) 的一个代表，第二条成立。

反过来，设第二条成立。取 \(\theta\in C^\infty(\mathbb R)\)，使 \(0\le\theta\le1\)，在 \((-\infty,1]\) 上为零，在 \([2,\infty)\) 上为一，并记 \(C=\|\theta'\|_\infty\)。对 \(0<\varepsilon<1/4\)，令
\[
 \chi_\varepsilon(x)=\theta(x/\varepsilon)
                 \theta\bigl((1-x)/\varepsilon\bigr).
\]
它在两端的 \(\varepsilon\) 邻域为零，在 \([2\varepsilon,1-2\varepsilon]\) 上为一，且
\(\|\chi_\varepsilon'\|_\infty\le2C/\varepsilon\)。
记
\[
 \omega_\varepsilon
 =\sup_{x\in[0,2\varepsilon]\cup[1-2\varepsilon,1]}|U'(x)|.
\]
由 \(U'\) 连续且端点值为零，\(\omega_\varepsilon\to0\)。利用 \(U\) 的两个端点值也为零，在这些边界区间上有
\[
 |U(x)|\le2\varepsilon\omega_\varepsilon.
\]
于是
\[
 \|(1-\chi_\varepsilon)U\|_\infty
 \le2\varepsilon\omega_\varepsilon,
 \quad
 \|\bigl((1-\chi_\varepsilon)U\bigr)'\|_\infty
 \le(1+4C)\omega_\varepsilon.
\]
因此 \(\chi_\varepsilon U\to U\) 于 \(W^{1,\infty}(I)\)。

每个 \(\chi_\varepsilon U\) 是内部紧支集的 \(C^1\) 函数，补零后仍属于 \(C_c^1(\mathbb R)\)。固定 \(\varepsilon\)，再选择充分小的磨光尺度，使卷积支集仍紧含于 \(I\)。命题\ref{b1:prop:local-smoothing-approximation}的 \(C^1\) 部分保证函数及其一阶导数一致逼近。依次选择这两个参数，便得到 \(C_c^\infty(I)\) 中依 \(W^{1,\infty}\) 逼近 \(u\) 的序列。

这一步只对已有 \(C^1\) 代表的函数使用逐阶一致磨光，并未声称任意 \(W^{1,\infty}\) 函数都有这样的逼近。
\end{solution}

\begin{exercise}\label{b1:ex:22-w0-order-ideal}
设 \(1\le p<\infty\)，实值 \(u\in W_0^{1,p}(\Omega)\)、\(v\in W^{1,p}(\Omega)\)，且 \(u\ge0\)、\(|v|\le u\) 几乎处处。
证明 \(v\in W_0^{1,p}(\Omega)\)。再在 \(I=(0,1)\) 上构造满足相同大小关系的 \(u\in C_c^\infty(I)\) 和内部紧支集 \(v\in W^{1,\infty}(I)\)，使 \(v\notin W_0^{1,\infty}(I)\)。
\end{exercise}
% 来源：自拟；用第21章正部强连续性制造保序逼近，无穷端点反例是内部尖点，不重复正文边界障碍。
\begin{solution}
取实值 \(\varphi_n\in C_c^\infty(\Omega)\)，使 \(\varphi_n\to u\) 于 \(W^{1,p}\)，并令 \(s_n=\varphi_n^+\)。由\cref{b1:prop:sobolev-lattice-strong-continuity}，\(s_n\to u^+=u\) 于同一空间。再令
\[
 v_n=\max\bigl(-s_n,\min(v,s_n)\bigr).
\]
正部运算、加法和减法都在有限指数的 \(W^{1,p}\) 中连续，而
\[
 \min(a,b)=a-(a-b)^+,\quad
 \max(a,b)=b+(a-b)^+.
\]
因此 \(v_n\in W^{1,p}(\Omega)\)，并在这个空间中有
\[
 v_n\longrightarrow\max\bigl(-u,\min(v,u)\bigr)=v.
\]
最后一个等式使用了 \(|v|\le u\)。

由于 \(-s_n\le v_n\le s_n\)，函数 \(v_n\) 在紧集 \(\operatorname{supp}\varphi_n\Subset\Omega\) 外几乎处处为零。它虽不一定光滑，却具有正文所定义的内部紧支集。由\cref{b1:prop:compact-sobolev-smooth-approximation}，每个 \(v_n\) 都能由内部紧支集光滑函数在 \(W^{1,p}\) 中逼近，故 \(v_n\in W_0^{1,p}(\Omega)\)。闭包空间本身是闭的，再由 \(v_n\to v\) 得到所求结论。

对无穷端点，取非负 \(\chi\in C_c^\infty(I)\)，使它在 \(1/2\) 的一个邻域等于一，令
\[
 u=\chi,\quad v(x)=\chi(x)|x-1/2|.
\]
则 \(0\le v\le u/2\le u\)。函数 \(v\) 为内部紧支集 Lipschitz 函数，因而属于 \(W^{1,\infty}(I)\)；但它在 \(1/2\) 处有尖点，不属于 \(C^1(I)\)。若它另有一个 \(C^1\) 代表，该代表与连续的 \(v\) 几乎处处相等，便由连续性处处相等，仍无法消除尖点。由\cref{b1:ex:22-w0-infty-interval}，\(v\notin W_0^{1,\infty}(I)\)。

所以内部紧支集本身不能保证无穷指数的光滑范数逼近；有限指数证明所用的非线性强连续性与内部磨光都确有指数范围。
\end{solution}

\begin{exercise}\label{b1:ex:22-halfspace-zero-extension-converse}
令 \(H=\{\,x\in\mathbb R^d\mid x_d>0\,\}\)，\(1\le p<\infty\)，并设 \(u\in W^{1,p}(H)\)。记 \(\widetilde u\) 为它在 \(H\) 外的零延拓。证明下列两条等价：
\begin{equivalences}
\item\label{b1:item:22-halfspace-w0}
\(u\in W_0^{1,p}(H)\)。
\item\label{b1:item:22-halfspace-global-zero-extension}
\(\widetilde u\in W^{1,p}(\mathbb R^d)\)。
\end{equivalences}
反向证明须实际构造支集留在半空间内部的光滑逼近，而不使用迹定理。
\end{exercise}
% 来源：自拟；由向内平移、远处截断和内部磨光证明半空间零延拓逆命题，正文只给正向与任意域的失败边界。
\begin{solution}
第一条蕴涵第二条是命题\ref{b1:prop:w0-zero-extension}。现设第二条成立，记 \(F=\widetilde u\)。对 \(h>0\) 定义
\[
 F_h(x)=F(x-he_d).
\]
这个函数在 \(\{x_d<h\}\) 上几乎处处为零。通过弱导数测试中的平移换元可得
\[
 \partial_jF_h(x)=(\partial_jF)(x-he_d).
\]
对函数和全部一阶导数分别应用\cref{b1:lem:lp-translation-continuity}，得到
\[
 \|F_h-F\|_{W^{1,p}(\mathbb R^d)}\longrightarrow0
 \quad(h\downarrow0).
\]
因而可以先把函数的支集与边界隔开，再处理它在无穷远处的部分。

固定 \(h\)。取 \(0\le\chi\le1\)、\(\chi\in C_c^\infty(\mathbb R^d)\)，使 \(\chi=1\) 于单位球，置 \(\chi_R(x)=\chi(x/R)\)、\(G_{h,R}=\chi_RF_h\)。当 \(R\to\infty\) 时，函数及导数中的
\[
 (\chi_R-1)F_h,\quad(\chi_R-1)\partial_jF_h
\]
都由控制收敛在 \(L^p\) 中趋零，而额外项满足
\[
 \|(\partial_j\chi_R)F_h\|_p
 \le R^{-1}\|\partial_j\chi\|_\infty\cdot\|F_h\|_p
 \longrightarrow0.
\]
所以 \(G_{h,R}\to F_h\) 于 \(W^{1,p}(\mathbb R^d)\)。同时 \(G_{h,R}\) 有紧支集，并在 \(\{x_d<h\}\) 上几乎处处为零。

最后固定 \(h,R\)，令
\[
 \varphi_{h,R,\varepsilon}=G_{h,R}*\rho_\varepsilon,
 \quad 0<\varepsilon<h/2.
\]
其支集有界且包含在 \(\{x_d\ge h-\varepsilon\}\) 中，因此 \(\varphi_{h,R,\varepsilon}\in C_c^\infty(H)\)。卷积微分公式以及有限指数的逼近恒等分别作用于 \(G_{h,R}\) 和各一阶导数，给出它在全空间 \(W^{1,p}\) 中的收敛：
\[
 \varphi_{h,R,\varepsilon}\longrightarrow G_{h,R}.
\]
给定任意误差要求，依次选择充分小的 \(h\)、充分大的 \(R\)、再选小于 \(h/2\) 的充分小 \(\varepsilon\)，使上述三段误差之和满足要求。限制到 \(H\) 不会增大 Sobolev 范数，所得函数便在 \(W^{1,p}(H)\) 中逼近 \(u\)。这证明 \(u\in W_0^{1,p}(H)\)。

半空间的作用在第一步：统一向内平移能为整个支集留下正距离。这个几何步骤在一般开集上未必可行，因此不能据本题撤去正文任意域逆命题的限制。
\end{solution}

\begin{exercise}[\(\star\) 选读]\label{b1:ex:22-second-order-reflection}
设 \(1\le p\le\infty\)，记 \(I=(0,\infty)\)。对 \(u\in W^{2,p}(I)\)，定义
\[
 Eu(x)=
 \begin{cases}
 u(x),&x>0,\\
 3u(-x)-2u(-2x),&x<0.
 \end{cases}
\]
原点处的值不影响等价类。证明 \(E\) 是从 \(W^{2,p}(I)\) 到 \(W^{2,p}(\mathbb R)\) 的有界线性延拓算子，并说明系数 \(3,-2\) 如何由原点处的两项匹配条件确定。
\end{exercise}
% 来源：自拟；修正正文偶反射的一阶构造，在半直线上显式匹配函数及一阶导数，不预借高阶迹理论。
\begin{solution}
先说明匹配条件有意义。固定 \(R>0\)，由有限区间上的 Hölder 不等式，\(u,u',u''\) 都在 \((0,R)\) 上可积。对 \(u'\) 应用\cref{b1:thm:one-dimensional-weak-ac}，其连续代表可写为某个常数加上 \(u''\) 的积分。由于 \(u''\in L^1(0,R)\)，这个代表绝对连续地延伸到原点，记为 \(G\)。同样，\(u\) 有一个代表 \(U\)，满足
\[
 U(x)=U(x_0)+\int_{x_0}^xG(t)\,\mathrm dt,
\]
并延伸为 \([0,R]\) 上的 \(C^1\) 函数，\(U'=G\)。在不同 \(R\) 上，连续代表由唯一性一致，故 \(U(0)\)、\(G(0)\) 都有明确含义。

假设先在负半轴使用 \(AU(-x)+BU(-2x)\)。其原点左极限为 \((A+B)U(0)\)，一阶导数的左极限为 \((-A-2B)G(0)\)。要对任意 \(u\) 分别与右侧的函数和导数匹配，应有
\[
 A+B=1,\quad -A-2B=1.
\]
解得 \(A=3\)、\(B=-2\)。因此题中延拓函数在原点连续，其分段一阶导数
\[
 Q(x)=
 \begin{cases}
 G(x),&x>0,\\
 -3G(-x)+4G(-2x),&x<0
 \end{cases}
\]
也在原点连续。

两侧绝对连续函数只要端点值相同，就能绝对连续地拼合：把两侧的积分表示以共同端点值为常数连接即可。由此，\(EU\) 在每个紧区间上绝对连续，弱导数是 \(Q\)；而 \(Q\) 同样绝对连续，其弱导数在负半轴上为 \(3u''(-x)-8u''(-2x)\)，在正半轴上为 \(u''(x)\)。两次拼合均无跳跃，因而没有遗漏集中于原点的分布项。

统一地，对 \(j=0,1,2\)，负半轴上的第 \(j\) 阶导数为
\[
 \partial^j(Eu)(x)
 =3(-1)^ju^{(j)}(-x)-2(-2)^ju^{(j)}(-2x).
\]
记 \(I_-=(-\infty,0)\)。换元给出
\[
 \|u^{(j)}(-2\,\cdot)\|_{L^p(I_-)}
 =2^{-1/p}\|u^{(j)}\|_{L^p(I)},
\]
其中 \(p=\infty\) 时约定 \(1/p=0\)。由三角不等式，负半轴上的导数范数至多是
\(\bigl(3+2^{j+1-1/p}\bigr)\|u^{(j)}\|_{L^p(I)}\)。
再加上正半轴的贡献，可使用统一粗界
\[
 \|\partial^j(Eu)\|_{L^p(\mathbb R)}
 \le12\|u^{(j)}\|_{L^p(I)}.
\]
对三个导数分量取相应的有限 \(p\) 和范数或无穷端点最大值，得到
\[
 \|Eu\|_{W^{2,p}(\mathbb R)}
 \le12\|u\|_{W^{2,p}(I)}.
\]
定义显然对等价类良定义并线性，且在正半轴保留 \(u\)，因此确为所需延拓算子。这里建立的是半直线上的二阶构造，不是任意 Lipschitz 开集的高阶延拓结论。
\end{solution}

\begin{exercise}\label{b1:ex:22-quadratic-localization}
设 \(\Omega\subset\mathbb R^d\) 为开集，实光滑函数 \(\eta_1,\ldots,\eta_N\) 满足
\[
 \sum_{j=1}^N\eta_j^2=1,\quad
 A=\left\|\sum_{j=1}^N|\nabla\eta_j|^2\right\|_{L^\infty(\Omega)}<\infty.
\]
证明对实值或复值 \(u\in H^1(\Omega)\)，每个 \(\eta_ju\) 都属于 \(H^1(\Omega)\)，并有
\[
 \sum_{j=1}^N\|\eta_ju\|_{H^1(\Omega)}^2
 =\|u\|_{H^1(\Omega)}^2
   +\int_\Omega\left(\sum_{j=1}^N|\nabla\eta_j|^2\right)|u|^2\,\mathrm dx.
\]
据此给出局部化范数和原范数的双向估计，并用一个具体例子说明右侧的额外项不能一般删除。
\end{exercise}
% 来源：自拟；平方和分解使交叉项精确抵消，新增可计算的局部化能量误差，不另命名未经核定的专门公式。
\begin{solution}
由平方和为一，\(|\eta_j|\le1\)；由 \(A<\infty\)，每个光滑 \(\nabla\eta_j\) 也有界。光滑乘法的弱 Leibniz 公式给出
\[
 \nabla(\eta_ju)=\eta_j\nabla u+u\nabla\eta_j.
\]
右侧两项均属于 \(L^2(\Omega)\)，而 \(\eta_ju\in L^2(\Omega)\)，所以乘积属于 \(H^1\)。这里不要求 \(u\) 本性有界：起作用的是光滑乘子的函数和一阶导数有界。

几乎处处展开欧氏长度的平方，
\[
 |\nabla(\eta_ju)|^2
 =\eta_j^2|\nabla u|^2+|u|^2\cdot|\nabla\eta_j|^2
   +2\eta_j\operatorname{Re}(\overline u\nabla u)\cdot\nabla\eta_j.
\]
对于实值函数，实部与共轭不改变表达式。对恒等式 \(\sum_j\eta_j^2=1\) 求导，得到
\[
 \sum_{j=1}^N\eta_j\nabla\eta_j=0.
\]
因此将前一等式求和时，全部交叉项相消，得到
\[
 \sum_{j=1}^N|\nabla(\eta_ju)|^2
 =|\nabla u|^2+\left(\sum_{j=1}^N|\nabla\eta_j|^2\right)|u|^2.
\]
零阶部分另有 \(\sum_j|\eta_ju|^2=|u|^2\)。将两条等式积分，便得所求能量恒等式。

额外项非负且至多为 \(A\|u\|_2^2\)，所以
\[
 \|u\|_{H^1}
 \le\left(\sum_{j=1}^N\|\eta_ju\|_{H^1}^2\right)^{1/2}
 \le(1+A)^{1/2}\|u\|_{H^1}.
\]
取 \(\Omega=(0,1)\)、\(\eta_1(x)=\cos x\)、\(\eta_2(x)=\sin x\)，以及 \(u=1\)。此时平方和为一，梯度平方和也为一，额外项恰等于 \(1\)，严格为正。局部化保持零阶平方和，却会为分割函数的变化付出一阶能量。
\end{solution}

\begin{exercise}\label{b1:ex:22-caccioppoli-reaction}
设实值 \(u\in H^1_{\mathrm{loc}}(\Omega)\)、\(f\in L^2_{\mathrm{loc}}(\Omega)\)，常数 \(\lambda>0\)，且对每个实值 \(\varphi\in C_c^\infty(\Omega)\) 都有
\[
 \int_\Omega\nabla u\cdot\nabla\varphi\,\mathrm dx
 +\lambda\int_\Omega u\varphi\,\mathrm dx
 =\int_\Omega f\varphi\,\mathrm dx.
\]
证明对每个非负 \(\eta\in C_c^\infty(\Omega)\)，
\[
 \int_\Omega\eta^2|\nabla u|^2\,\mathrm dx
 +\lambda\int_\Omega\eta^2u^2\,\mathrm dx
 \le4\int_\Omega u^2|\nabla\eta|^2\,\mathrm dx
     +\lambda^{-1}\int_\Omega\eta^2f^2\,\mathrm dx.
\]
须先说明为什么可以用通常不光滑的 \(\eta^2u\) 代入测试恒等式。
\end{exercise}
% 来源：自拟；把正文弱调和Caccioppoli方法扩展到源项及正零阶项，独立说明测试合法化及吸收常数。
\begin{solution}
令 \(\psi=\eta^2u\)。光滑乘法的弱求导公式说明 \(\psi\in H^1(\Omega)\)，且它具有内部紧支集。由\cref{b1:prop:compact-sobolev-smooth-approximation}，存在 \(\psi_n\in C_c^\infty(\Omega)\) 在 \(H^1\) 中逼近它。还可以要求这些测试函数的支集包含在同一个开集 \(V\) 中，并满足 \(\overline V\Subset\Omega\)：先选 \(\zeta\in C_c^\infty(V)\)，使 \(\zeta=1\) 于 \(\operatorname{supp}\eta\) 的邻域，再把任意所述逼近列换成 \(\zeta\psi_n\)。乘子及其梯度有界，所以替换后仍趋于 \(\zeta\psi=\psi\)。

函数 \(u,\nabla u,f\) 在 \(V\) 上均平方可积。测试恒等式三项之差的绝对值分别受
\[
 \begin{gathered}
 \|\nabla u\|_{L^2(V)}\cdot\|\nabla(\psi_n-\psi)\|_{L^2(V)},\\
 \lambda\|u\|_{L^2(V)}\cdot\|\psi_n-\psi\|_{L^2(V)},\\
 \|f\|_{L^2(V)}\cdot\|\psi_n-\psi\|_{L^2(V)}
 \end{gathered}
\]
控制，都趋于零。因而恒等式对 \(\psi=\eta^2u\) 成立。

利用 \(\nabla(\eta^2u)=\eta^2\nabla u+2\eta u\nabla\eta\)，代入得
\[
 \int_\Omega\eta^2|\nabla u|^2
 +\lambda\int_\Omega\eta^2u^2
 =-2\int_\Omega\eta u\nabla u\cdot\nabla\eta
   +\int_\Omega\eta^2fu.
\]
分别使用两个带参数的平方不等式，
\[
 2|\eta|\cdot|u|\cdot|\nabla u|\cdot|\nabla\eta|
 \le\frac12\eta^2|\nabla u|^2+2u^2|\nabla\eta|^2,
\]
\[
 \eta^2|f|\cdot|u|
 \le\frac{\lambda}{2}\eta^2u^2+\frac1{2\lambda}\eta^2f^2.
\]
将它们积分并代回，梯度项和零阶项各有一半可以吸收到左端；再将两边乘以 \(2\)，恰得题中的估计。常数 \(\lambda>0\) 在第二个吸收步骤中使用，不能在最终公式里直接令 \(\lambda=0\)。
\end{solution}

\begin{exercise}[\(\star\) 选读]\label{b1:ex:22-dirichlet-minimum}
设 \(\Omega\subset\mathbb R^d\) 为非空有界开集，实值 \(g\in H^1(\Omega)\)，记
\[
 \mathcal A=g+H_0^1(\Omega),\quad
 \mathcal E(v)=\int_\Omega|\nabla v|^2\,\mathrm dx.
\]
证明对 \(u\in\mathcal A\)，下列两条等价：
\begin{equivalences}
\item\label{b1:item:22-energy-minimum}
\(\mathcal E(u)=\inf_{v\in\mathcal A}\mathcal E(v)\)。
\item\label{b1:item:22-energy-orthogonality}
对每个实值 \(\varphi\in H_0^1(\Omega)\)，都有
\[
 \int_\Omega\nabla u\cdot\nabla\varphi\,\mathrm dx=0.
\]
\end{equivalences}
再证明这样的 \(u\) 存在且唯一，并说明这里的边界条件为何不需要先定义边界迹。
\end{exercise}
% 来源：自拟；题1恢复纯梯度内积的正定性，再用第8章Hilbert Riesz求仿射闭包类的最小能量元素。
\begin{solution}
先证明等价性。若 \(u\) 为极小点，对任意 \(\varphi\in H_0^1(\Omega)\) 和实数 \(t\)，仍有 \(u+t\varphi\in\mathcal A\)，而
\[
 \mathcal E(u+t\varphi)
 =\mathcal E(u)+2t\int_\Omega\nabla u\cdot\nabla\varphi
                 +t^2\|\nabla\varphi\|_2^2.
\]
这个关于 \(t\) 的二次多项式在 \(t=0\) 取极小值，其一次项系数必须为零，所以第二条成立。

反过来，若第二条成立，任意 \(v\in\mathcal A\) 都可写成 \(v=u+\psi\)，其中 \(\psi\in H_0^1(\Omega)\)。交叉项为零，故
\[
 \mathcal E(v)=\mathcal E(u)+\|\nabla(v-u)\|_2^2
 \ge\mathcal E(u).
\]
这就证明了第一条。

为证明存在性，在 \(V=H_0^1(\Omega)\) 上定义实内积
\[
 (v,w)_V=\int_\Omega\nabla v\cdot\nabla w\,\mathrm dx.
\]
因为 \(\Omega\) 有界，可以把它放入某个有限宽度的坐标条带。习题\ref{b1:ex:22-strip-poincare}说明 \(\|\nabla v\|_2\) 在 \(V\) 上确为范数，并与 \(H^1\) 范数等价。空间 \(V\) 按定义为 \(H^1(\Omega)\) 的闭子空间，故关于原范数完备；范数等价使其关于梯度内积也完备。因此这是一个 Hilbert 空间。

定义连续线性泛函
\[
 \Lambda(\varphi)=-\int_\Omega\nabla g\cdot\nabla\varphi\,\mathrm dx.
\]
由 Cauchy--Schwarz 不等式，
\[
 |\Lambda(\varphi)|\le\|\nabla g\|_2\cdot\|\varphi\|_V.
\]
Hilbert 空间的 Riesz 表示定理\ref{b1:thm:hilbert-riesz}给出 \(h\in V\)，使
\[
 \int_\Omega\nabla h\cdot\nabla\varphi\,\mathrm dx
 =-\int_\Omega\nabla g\cdot\nabla\varphi\,\mathrm dx
 \quad(\varphi\in V).
\]
令 \(u=g+h\)，则 \(u\in\mathcal A\) 且满足正交条件，因而是极小点。

若 \(u_1,u_2\) 都为极小点，上面的能量分解给出
\(\|\nabla(u_1-u_2)\|_2=0\)。其差属于 \(H_0^1(\Omega)\)，再次使用题1的不等式便得 \(u_1=u_2\) 几乎处处。这证明在 Sobolev 等价类意义下唯一；连通性并不是这里所需的额外假设。

本题的边界条件就是 \(u-g\in H_0^1(\Omega)\)，即这个差能由内部紧支集光滑函数依 \(H^1\) 范数逼近。整个构造没有给每个边界点赋值，也没有将任意开集上的闭包条件认作某个尚未建立的零迹条件。若 \(g\) 换成同一仿射类的另一个代表，集合 \(\mathcal A\) 不变，所得唯一极小点也不变。
\end{solution}
